% Modernised reading — fonds Alexandre Grothendieck, shelfmark 156-9, pages 1-139.
%
% This is an INTERPRETATION, not a transcription. Everything the critical
% apparatus carried moves into footnotes. In any dispute of substance the
% transcription governs: whoever cites Grothendieck must cite batch-NN.fr.tex.
%
% Pass: Opus 5.5 (claude-opus-5-5), 2026-10-10 — first pass, unchecked against
% the pages by a human. The folder was read whole: seven batches, pages 1-139,
% dated 5 July (p. 4), 8 July (struck, p. 33), 10 July (p. 75), 11 July
% (doubtful, p. 98) and 15 July 1986 (pp. 2-3, written last and filed first).
% Pages 109, 111, 113 are typed versos of another text and 128 is blank; none
% is transcribed or restated. Not measured against the Opus 5 reference
% reading of folder 115.
%
% What the folder is. Chapter IX (« GF IX ») of Vers une géométrie des formes,
% and the opening of IX bis (p. 129 on, his new pagination 1-11). It follows
% the fourth mouture of Analysis situs (GF VIII = 156-8), whose axioms Mag 1-4
% it takes for granted. Spine: (1) the atelier of an ordered set (a
% combinatorial segment), its figures (finite chains) and the duality
% figure / complementary figure, pp. 4-27; (2) four digressions, the last of
% which swaps « atelier » and « magasin », pp. 28-40; (3) supports, spatial
% ateliers and « spatios » (complete ortholattices), pp. 40-60; (4)
% dispositions (orthogonality spaces) and their correspondences — a first
% attempt pp. 60-74, restarted on 10 July pp. 75-100 — ending in a self-dual
% category; (5) set-theoretic analogy, orthomodularity, admissible families
% and the support axioms Atmagsupp 1-3, pp. 101-121; (6) a philosophical
% critique and a programme, pp. 121-127; (7) IX bis: morphisms of ateliers,
% excision, subdivisions, pp. 129-139, where the folder stops.
%
% Notation fixed for the whole document. « Atelier » 𝒜 = set of strata with ⊴
% (sub-stratum), ◁ (proper face), ≪, °≪ (\mathring{\ll}), |o| (disjunction)
% — the terminology he adopts on p. 40; pp. 4-39 say « magasin » and write 𝔐
% for it, and « atelier » for the set of figures, which we call the magasin
% (attenant) 𝓕. L an ordered set; 𝒜(L), 𝒜̂(L); < the « left of » order on
% strata; δX, ∂X, |X|, |X|°; or, ex. Préfigure = his « polygone » of pp.
% 21-32. Drapeau = finite chain (his « diagramme » on pp. 21-26). Σ_𝒜, supp°,
% cosupp°, X°, ∁, Omb°. Spatios (Σ, ≤, ∁) for what pp. 77-80 call
% « spatiale » and pp. 81-99 « système ». Disposition (E, |o|) — « disposante »
% on pp. 101-102. supp° where pp. 81-99 write Sup°, Cosup°. α, ᾱ, α_*, β.
% ≍ for the compatibility relation of chapter VIII (p. 6; \between on pp.
% 129-130): our identification.
%
% Substantive departures from the pages, each footnoted where it occurs:
% — p. 2: the « modular » law written is infinite distributivity; in a spatios
%   it makes Σ a complete Boolean algebra.
% — p. 9: X °≪ I_{<a} iff X < a (the page writes X ≤ a; X = a is incident).
% — p. 11: the lemma is « il faut et il suffit »; the page proves the
%   sufficiency for Y in 𝒜 only.
% — p. 13: redundant vertex defined by « all neighbouring strata in 𝔉 »
%   (ord = 2 fails at an extreme of L); 𝔉** described with °≪, not ≪.
% — p. 15: the theorem is given a proof (none on the page).
% — pp. 7, 19, 20: « contractibles » read as « constructibles ».
% — p. 20: 𝔉_is is the set of ≪-maximal (not °≪-maximal) vertices of Σ.
% — p. 22: injectivity of Φ ↦ S_Φ supplied (the count 2^{2n+1} needs it).
% — p. 27: the « resp. » of the divisibility conditions is inverted.
% — p. 32: the claim that {X : X |o| Φ, X ⋡ Φ} is a sub-atelier is false in
%   dimension 2 (counter-example given); the page's proof treats only Y' ∈ Φ;
%   IX bis (p. 130) states the correct version, with compatibility.
% — p. 49: Atspat 2 needs X ≠ Y.
% — p. 71: « Lemme » (bijection Hom(E, Σ_E') ≅ sup-maps) false as stated;
%   only the injection holds, its image is Prop. 4's condition.
% — p. 84: in (v), x', y' range over E'; (ii) stated as in Cor. 2, b').
% — p. 95: (iv) ⇔ (v) of Prop. 4 proved directly; the page's added clause
%   is not needed.
% — p. 100: the product relation is (x,x') |o| (y,y') iff x |o| y or x' |o| y'.
% — p. 105: commuting with ∁ means f' bijective (page: f); g = f∘f'^{-1}.
% — p. 107: the Question has a positive answer, proved in a footnote.
% — p. 110: the doubts are right: O6 (hexagon) is a finite non-orthomodular
%   ortholattice, MO2 orthomodular and not Boolean.
% — p. 118: A and B are swapped in Atmagsupp 2; p. 120: the union runs from 0.
% — p. 122: the failure of transitivity needs Σ non-orthomodular, not merely
%   non-Boolean (Foulis–Holland).
%
% Modern names given and footnoted as ours: relation de polarité et connexion
% de Galois (Birkhoff 1940, Ore 1944), espace d'orthogonalité (Dacey 1968,
% Foulis–Randall), treillis orthocomplémenté complet, complétion de
% Dedekind–MacNeille (1937), sup-treillis et adjoints, catégorie à poignard
% (Jacobs 2010), treillis orthomodulaire (Husimi 1937, Kalmbach 1983),
% théorème de Foulis–Holland, O6 et MO2, ensembles constructibles de la droite
% et structures o-minimales (van den Dries 1984, Pillay–Steinhorn 1986),
% produit tensoriel en logique quantique (Foulis–Randall 1981).
%
% What the folder announces and never establishes: Cor. 2 of p. 15
% (unfinished); the case of two infinite intervals in the lemma of p. 11; the
% claim of p. 57 that Atspat 1', 2' are « les deux seules (?) » conditions,
% beyond what p. 59 checks; the end of Lemma 3 (p. 120); whether Atmagsupp 3
% follows from a faithful support function (p. 123); the structure (∁, |o|) on
% Corr(Σ, Σ') and the open questions of pp. 99-100; the decomposition of a
% general correspondence into comorphism and morphism (p. 103); the excision
% 𝒜/N (pp. 134-136); (iii) ⇒ (i) for subdivisions (p. 138); the programme of
% p. 127 (subdivisions, moderate ateliers, saturation); whether moderate
% ateliers are spatial and faithful (p. 3).
\documentclass[11pt,a4paper]{article}
\input{../preamble/grothendieck}

\folder{156-9}
\pages{1}{139}
\dating{1986}
\watermark{Édition de démonstration\\interprétation personnelle de l'œuvre}
\foldertitle{[Chapitre] IX et IX bis. [Ateliers] : notes manuscrites (05-15/07/1986) — lecture modernisée du dossier entier}

\begin{document}

\section*{Résumé}

\begin{resume}
On peut décrire une forme géométrique sans parler de ses points : un triangle,
c'est trois sommets, trois arêtes ouvertes et un intérieur ouvert, avec trois
renseignements --- qui est la face de qui, qui est contenu dans qui, qui est
disjoint de qui. Grothendieck appelle \emph{strates} ces morceaux, et
\emph{atelier} l'ensemble des strates muni de ces trois relations ; les
\emph{figures} sont les assemblages finis de strates, comme les subdivisions
d'un polyèdre. Ces cent trente-neuf feuillets, écrits du 5 au 15 juillet 1986,
forment le chapitre IX de \emph{Vers une géométrie des formes} et le début d'un
chapitre IX bis ; ils viennent juste après la quatrième mouture d'« Analysis
situs » (chapitre VIII), qui avait posé les axiomes.

Le dossier commence par l'exemple le plus simple : le segment. Un ensemble
ordonné $L$ donne des strates --- ses points, et les intervalles ouverts entre
deux points --- et une figure n'est rien d'autre qu'une suite finie de points et
d'intervalles rangés de gauche à droite. Le premier résultat est une dualité :
ce qui, dans le segment, est disjoint d'une figure est encore décrit par une
figure, la \emph{figure complémentaire}, à condition d'admettre des intervalles
qui partent à l'infini. Qu'on pense à un ruban marqué de quelques traits :
découper les morceaux qu'on garde laisse des chutes, et les chutes sont encore
des traits et des morceaux.

Vient ensuite la question centrale : qu'est-ce que la \emph{place occupée} par
une strate, si l'on ne dispose d'aucun point ? La réponse ne se sert que de la
disjonction. On appelle \emph{support} d'un ensemble de strates tout ce qui est
disjoint de tout ce qui lui est disjoint --- l'« orthogonal de l'orthogonal »,
comme en géométrie euclidienne l'orthogonal de l'orthogonal d'une droite est la
droite elle-même. Les supports s'ordonnent par inclusion et ont un complément.
Grothendieck isole les propriétés de cette structure sous le nom de
\emph{spatios}, et découvre que, pour les ateliers qu'il dit « spatiaux », on
peut tout reconstruire à partir des supports : l'ensemble de base de la théorie
change, une fois de plus, et devient le « spatios des supports ». Il étudie
ensuite les applications entre ces structures, constate que chacune a une
\emph{transposée}, et note lui-même l'analogie avec les opérateurs adjoints des
espaces de Hilbert.

Cette analogie n'est pas fortuite. Les spatios sont ce qu'on appelle
aujourd'hui des treillis orthocomplémentés complets, et la condition qu'il
cherche ensuite pour ses « axiomes des supports » --- que les supports des
strates d'une figure se comportent comme des parties d'un ensemble --- revient
exactement à l'\emph{orthomodularité}, la loi qui distingue les sous-espaces
fermés d'un espace de Hilbert des treillis quelconques. Il pressent qu'elle peut
manquer (« J'en doute fort ! »), et il a raison : un hexagone en est le
contre-exemple le plus petit. Le dossier aboutit ainsi, venu de la géométrie,
au cadre de la logique quantique, sans le nommer.

Ces pages montrent aussi un travail qui se reprend sans cesse. Le 5 juillet :
« L'affaire a tendance à s'enliser, et j'ai envie de tout reprendre » ; le 10 :
« Je reprends dès le début ». Une page avoue : « je n'ai pas l'impression encore
d'avoir vraiment bien compris la situation » ; une autre conclut une
démonstration par « On gagne ! ». Les noms changent sous nos yeux --- atelier et
magasin échangent leurs rôles --- parce que, dit-il, l'accent s'est déplacé. Et
une section s'intitule « Critique philosophique de cette approche » : la
construction l'oblige à découper des strates qui ne se comparent pas, alors
qu'il ne s'est jamais intéressé qu'aux supports constructibles ; il y voit une
contradiction, et note que la « spatialité » a « de forts relents de topologie
générale ». Le dossier s'achève sur les morphismes d'ateliers et une définition
des subdivisions, interrompue.

\keywords{interval poset, combinatorial segment, constructible set, tame topology, orthogonality space, Galois connection, closure operator, complete ortholattice, join-dense subset, Dedekind–MacNeille completion, sup-lattice, adjoint map, dagger category, orthomodular lattice, commuting elements, locally finite family, sub-poset closed under faces, subdivision}
\end{resume}

\section*{Le fil du dossier, et les conventions}

\begin{enumerate}
\item[1.] \emph{Le tableau des types d'ateliers} (pages 1 à 3), daté du 15
juillet : écrit en dernier, placé en tête.
\item[2.] \emph{L'atelier d'un ensemble ordonné} (pages 4 à 6), 5 juillet :
strates, relations, et le fait que les préfigures sont les chaînes finies.
\item[3.] \emph{Intervalles infinis et préfigure complémentaire} (pages 7 à
15), jusqu'au théorème $\operatorname{cosupp}^{\circ}(\mathfrak{F}) =
\operatorname{Omb}^{\circ}(\mathfrak{F}^{*})$.
\item[4.] \emph{Autre présentation : le segment orienté combinatoire} (pages 16
à 20), seconde rédaction du même sujet, avec la classification des sommets et
leur échange par dualité.
\item[5.] \emph{Supports subordonnés à un drapeau ; segments ouverts} (pages 21
à 27).
\item[6.] \emph{Quatre digressions} (pages 28 à 40) : somme, sous-ateliers,
compactologie, terminologie.
\item[7.] \emph{Supports, ateliers spatiaux, spatios} (pages 40 à 60).
\item[8.] \emph{Premier essai sur les correspondances de dispositions} (pages 60
à 74).
\item[9.] \emph{Reprise du 10 juillet : dispositions et spatios} (pages 75 à
87).
\item[10.] \emph{Correspondances et transposition} (pages 87 à 100).
\item[11.] \emph{L'analogie ensembliste et ses questions} (pages 101 à 107).
\item[12.] \emph{Orthomodularité, familles admissibles, axiomes des supports}
(pages 108 à 121).
\item[13.] \emph{Critique philosophique et programme} (pages 121 à 127).
\item[14.] \emph{IX bis : morphismes d'ateliers, excision, subdivisions} (pages
129 à 139).
\end{enumerate}

Les feuillets 4 à 127 portent sa pagination 1 à 122 (la page 4 des archivistes
est sa page 1, avec quelques numéros récrits) ; une nouvelle pagination 1 à 11
commence à la page 129, sous l'en-tête « cf IX bis ». Les renvois de sa main
(« p.~59 ») sont à sa pagination, et nous les traduisons. Les pages 109, 111 et
113 sont des versos dactylographiés d'un autre texte, sans rapport avec ces
notes ; la page 128 est blanche. Deux sections sont des \emph{moutures}
successives d'un même sujet, et sont gardées distinctes : les pages 7 à 15 et
16 à 20 (deux présentations du segment), les pages 60 à 74 et 75 à 100 (deux
théories des correspondances).

\emph{Ce que le dossier suppose.} Les axiomes du chapitre VIII (dossier 156-8),
qu'il appelle Mag~1 à Mag~4 jusqu'à la page 40 et At~1 à At~4 ensuite. Un
atelier est un ensemble $\mathcal{A}$ muni de trois relations : $X
\trianglelefteq Y$ (« $X$ est une sous-strate, une face, de $Y$ »), $X \ll Y$
(« $X$ est contenu dans l'adhérence de $Y$ ») et $X \mathrel{|\circ|} Y$ (« $X$
et $Y$ sont disjoints », les strates étant vues comme ouvertes), avec : At~1,
$\trianglelefteq$ et $\ll$ sont des relations d'ordre et $\trianglelefteq$
entraîne $\ll$ ; At~2, si $X \ll Y$, l'ensemble des $Y' \trianglelefteq Y$ tels
que $X \ll Y'$ a un plus petit élément, et l'on écrit $X \mathrel{\mathring{\ll}}
Y$ quand c'est $Y$ lui-même (« $X$ est contenu dans la strate ouverte $Y$ ») ;
At~3, $\mathrel{|\circ|}$ est symétrique et antiréflexive, et $X \mathrel{|\circ|}
Y$, $X' \mathrel{\mathring{\ll}} X$, $Y' \mathrel{\mathring{\ll}} Y$ entraînent
$X' \mathrel{|\circ|} Y'$ ; At~4, deux sous-strates distinctes d'une même strate
sont disjointes. On note $X \lhd Y$ pour $X \trianglelefteq Y$, $X \neq Y$, et
$\widetilde{X} = \{X' \mid X' \trianglelefteq X\}$. Une \emph{figure} est une
partie finie, fermée pour $\trianglelefteq$, dont les éléments sont deux à deux
compatibles au sens du chapitre VIII ; une \emph{préfigure} est une partie dont
la fermeture pour $\trianglelefteq$ est une figure.

\emph{Conventions.} Nous disons partout \emph{atelier} pour l'ensemble des
strates, noté $\mathcal{A}$, et \emph{magasin} pour un ensemble de figures,
noté $\mathcal{F}$ : c'est la terminologie qu'il adopte à la page 40 et
qu'emploient déjà les pages 1 à 3, écrites en dernier. Les pages 4 à 39 disent
l'inverse et écrivent $\mathfrak{M}$ ; nous écrivons $\mathcal{A}(L)$ là où
elles écrivent $\mathfrak{M}$. Pour un ensemble ordonné $L$, $<$ désigne aussi
l'ordre « à gauche de » entre strates de $\mathcal{A}(L)$, qui le prolonge ; il
ne faut pas le confondre avec $\trianglelefteq$. Un \emph{drapeau} est une
partie finie totalement ordonnée (une chaîne) ; les pages 21 à 26 disent
« diagramme ». Nous disons \emph{préfigure} là où les pages 21 à 32 disent
« polygone ». Pour une partie $A$ de $\mathcal{A}$,
\[
\operatorname{cosupp}^{\circ}(A) = \{ Z \in \mathcal{A} \mid Z
\mathrel{|\circ|} X \ \forall X \in A \}, \qquad \operatorname{supp}^{\circ}(A)
= \operatorname{cosupp}^{\circ}\operatorname{cosupp}^{\circ}(A),
\]
$\Sigma_{\mathcal{A}}$ est l'ensemble des supports (les parties de la forme
$\operatorname{cosupp}^{\circ}(B)$), $X^{\circ} = \operatorname{supp}^{\circ}\{X\}$,
$\complement S = \operatorname{cosupp}^{\circ}(S)$ pour $S \in
\Sigma_{\mathcal{A}}$, et $\operatorname{Omb}^{\circ}(A) = \{Z \mid Z
\mathrel{\mathring{\ll}} X \text{ pour un } X \in A\}$ est l'\emph{ombre} de $A$.
Un \emph{spatios} $(\Sigma, \leq, \complement)$ est ce que les pages 77 à 80
appellent « spatiale » et les pages 81 à 99 « système ». Une \emph{disposition}
est un ensemble muni d'une relation symétrique $\mathrel{|\circ|}$ (« disposante »
aux pages 101 et 102). Les pages 81 à 99 écrivent $\operatorname{Sup}^{\circ}$ et
$\operatorname{Cosup}^{\circ}$ pour nos $\operatorname{supp}^{\circ}$ et
$\operatorname{cosupp}^{\circ}$.

\emph{Ce que le dossier annonce et n'établit pas} : le corollaire 2 de la page
15, inachevé ; le cas de deux intervalles infinis dans le lemme de la page 11 ;
la fin du lemme 3 (page 120) ; la question de savoir si une fonction support
fidèle entraîne l'axiome Atmagsupp~3 (page 123) ; une structure de complément
sur les correspondances entre spatios, et les « questions ouvertes » des pages
99 et 100 ; la décomposition d'une correspondance quelconque en comorphisme et
morphisme (page 103) ; l'excision $\mathcal{A}/N$ (pages 134 à 136) ;
l'implication (iii) $\Rightarrow$ (i) pour les subdivisions (page 138) ; le
programme de la page 127 (subdivisions, ateliers modérés, saturation) ; et la
question, posée en marge de la page 3, de savoir si un atelier modéré est
nécessairement spatial et fidèle.

\pagerange{1}{3}
\section*{Le tableau des types d'ateliers (pages 1 à 3)}

La page 1 porte en tête « GF IX », et au crayon un renvoi à « p.~75 [pour] la
mise au point sur la terminologie »\footnote{Le renvoi est à une page 75 de sa
main, que nous ne savons pas situer avec certitude ; la mise au point
terminologique de ce dossier est aux pages 37 à 40 (ses pages 34 à 37).}. Les
pages 1 et 3, datées du 15 juillet, sont un tableau des espèces d'ateliers
rencontrées jusque-là, avec leurs inclusions. La version propre, à la page 3,
est un graphe dont les flèches vont d'une classe vers une classe plus large :
\begin{itemize}
\item G, ateliers généraux ; au-dessous S (spatiaux), F (fidèles), L (locaux) ;
\item SF, spatiaux fidèles\footnote{La légende écrit « spatiaux locaux », ce qui
contredit le sigle et la place dans le graphe (SF s'envoie sur S et sur F) ; le
transcripteur signale déjà ce lapsus.} ; LS, locaux spatiaux, dont la légende
dit qu'ils sont les ateliers spatiaux divisibles ; P $=$ FL, fidèles locaux ou
« ponctuaires » ;
\item E $=$ FLS, fidèles locaux spatiaux, ou \emph{ensemblistes} ; M, modérés ;
EM, ensemblistes modérés ; Mq, maquettes ;
\item D, discrets ; et tout en bas p (ponctuels, un seul élément) et v (vide).
\end{itemize}
Les huit « carrés élémentaires » du graphe sont des carrés d'intersection --- la
classe du bas est l'intersection des deux classes du milieu --- et une « chasse »
dans le diagramme donne par exemple $\mathrm{LS} \cap \mathrm{M} = \mathrm{EM} =
\mathrm{E} \cap \mathrm{M}$. Les définitions de ces classes (fidèle, local,
ponctuaire, maquette) sont celles du chapitre VIII ; ce dossier ne les répète
pas\footnote{Au chapitre VIII (dossier 156-8, première page-tableau), un atelier
est ponctuaire quand $\Sigma \simeq \mathfrak{P}(L)$ pour un ensemble $L$ de
strates-points, une maquette quand $\Sigma \simeq \mathfrak{P}(\mathcal{A})$, et
fidèle quand l'ombre est une fonction support fidèle. La première version du
tableau, à la page 1, place aussi un « atelier divisible » sous l'atelier
général ; à la page 3 la divisibilité est un point marqué sur une arête.}.

Les cadres privilégiés pour une géométrie des formes lui semblent maintenant LS,
pour l'étude des « cordes et formes » (pseudo-simplexes topologiques), S pour
les contextes finis, M pour la topologie modérée, et EM pour la topologie
modérée ensembliste. Deux remarques en marge : il faudrait vérifier qu'un
atelier modéré est non seulement spatial mais fidèle, « sinon, diagramme moins
joli ! » ; et « l'accent principal s'est nettement déplacé depuis les approches
originelles via L, vers des approches subordonnées à S ». Le dossier entier est
l'histoire de ce déplacement : on verra, aux pages 40 à 60, ce que « spatial »
veut dire.

Les pages 1 et 2 ajoutent, dans une écriture très rapide\footnote{La lecture du
paragraphe de la page 1 est très incertaine dans la transcription ; on n'en
retient que ce qui est lisible.}, qu'il s'intéressera surtout aux ateliers
spatiaux, et parmi eux aux ateliers locaux et modérés. Il se demande s'il faut
exiger que le treillis $\Sigma$ des supports vérifie
\[
A \wedge \operatorname{Sup}_{i} B_{i} = \operatorname{Sup}_{i} (A \wedge B_{i}),
\]
et présume que cela entraînerait la fidélité\footnote{Il appelle cette loi
« modulaire ». C'est en fait la distributivité infinie de $\wedge$ sur les Sup,
plus forte que la modularité au sens de la théorie des treillis. Dans un spatios
(pages 49 et 77), elle fait de $\Sigma$ une algèbre de Boole complète : un
treillis orthocomplémenté distributif est booléen. C'est l'hypothèse
« booléenne » qui revient aux pages 106, 122 et 126. La marge écrit « À vérifier ?
Non », d'une lecture douteuse.}. Il pose enfin une question : pour une strate
ponctuelle $x$ d'un atelier local, le support $x^{\circ}$ est-il un élément
minimal de $\Sigma \smallsetminus \{0\}$ ? Autrement dit, si $X \neq x$,
existe-t-il $Z$ disjoint de $x$ et non disjoint de $X$ ? Il ramène la question
au cas où $X$ est lui-même un point $y$, et la laisse ouverte\footnote{La
réduction utilise $x^{\circ} = \operatorname{Sup} y^{\circ}$ sur les points $y$
de $X$ ; la phrase qui suit est en partie illisible. Une remarque datée du 15
juillet doute que la restriction d'un atelier à un « contexte » se fasse dans
les ateliers spatiaux sans plus ; un second graphe, à sigles, double le
tableau.}.

\pagerange{4}{6}
\section*{L'atelier d'un ensemble ordonné (pages 4 à 6)}

« 5 juillet. L'affaire a tendance à s'enliser, et j'ai envie de tout reprendre,
d'une façon un peu différente. » Ce qui suit reprend la fin du chapitre VIII,
où les ateliers attachés à un ensemble ordonné venaient d'être introduits comme
premier cas où les axiomes des supports ne sont pas évidents.

\subsection*{Strates et relations}

Soit $L$ un ensemble ordonné, pas nécessairement totalement ordonné. L'atelier
$\mathcal{A}(L) = \mathcal{A}_{0} \sqcup \mathcal{A}_{1}$ a pour strates :
\begin{itemize}
\item les \emph{sommets} ou \emph{lieux} $x \in \mathcal{A}_{0} = L$ ;
\item les \emph{intervalles} $I_{a,b}$ ($= I_{b,a}$), un pour chaque paire $a <
b$ de $L$, c'est-à-dire pour chaque drapeau à deux éléments.
\end{itemize}
On pose $\partial I_{a,b} = \delta I_{a,b} = \{a,b\}$, $\partial x = \emptyset$,
$\delta x = \{x\}$ ; le support dans $L$ est $|x| = \{x\}$, $|I_{a,b}| = [a,b]
= \{x \mid a \leq x \leq b\}$, et l'intérieur $|X|^{\circ} = |X|
\smallsetminus \partial X$. L'incidence est $x \lhd I_{a,b}$ si $x \in \{a,
b\}$, et $X \trianglelefteq Y$ veut dire $X \lhd Y$ ou $X = Y$. On pense à $L$
comme à un segment, à $x$ comme à un point et à $I_{a,b}$ comme à l'intervalle
ouvert $]a,b[$.

Les deux relations d'inclusion sont
\[
X \ll Y \iff |X| \subset |Y|, \qquad X \mathrel{\mathring{\ll}} Y \iff |X|
\subset |Y| \text{ et non } X \lhd Y,
\]
soit explicitement : $X \mathrel{\mathring{\ll}} Y$ si $X = Y$, ou $X = x$, $Y =
I_{a,b}$ avec $a < x < b$, ou $X = I_{a,b}$, $Y = I_{a',b'}$ avec $a' \leq a < b
\leq b'$ ; et $X \ll Y$ dans les mêmes cas, sauf que $a \leq x \leq b$ est
permis. La strate ouverte $X$ est contenue dans la strate ouverte $Y$, ou dans
son adhérence. On a $X \ll Y$ si et seulement si $X \mathrel{\mathring{\ll}} Y'
\trianglelefteq Y$ pour un $Y'$, et alors $X \lhd Y$ ou $X
\mathrel{\mathring{\ll}} Y$, les deux cas s'excluant. Enfin $X
\mathrel{\mathring{\ll}} Y$ entraîne que $\delta X \cup \delta Y$ est un drapeau.

\subsection*{L'ordre « à gauche de », et la disjonction}

Avant de définir la disjonction, on ordonne les strates :
\[
X < Y \iff X \neq Y \text{ et } x \leq y \text{ pour tous } x \in \delta X,\ y
\in \delta Y ,
\]
c'est-à-dire $x < y$ ; $x \leq a$ pour $Y = I_{a,b}$ ; $b \leq y$ pour $X =
I_{a,b}$ ; $b \leq a'$ pour deux intervalles. On a $a < I_{a,b} < b$. C'est un
ordre strict, qui prolonge celui de $L$. Avec $\operatorname{or} X$ et
$\operatorname{ex} X$, le plus petit et le plus grand élément de $\delta X$, les
remarques de la page 5 notent : a) $X < Y \Rightarrow \operatorname{ex} X \leq
\operatorname{or} Y$ ; b) si $X \leq Y$, alors $\operatorname{or} X <
\operatorname{ex} Y$, sauf si $X = Y$ est un point.

\begin{quote}
\textbf{Proposition} (page 5). \emph{Si $X < Y$, $X' \ll X$ et $Y' \ll Y$, alors
$X' < Y'$, ou bien $X' = \operatorname{ex} X = Y' = \operatorname{or} Y$.}

\textbf{Corollaire.} \emph{Si $X < Y$, $X' \mathrel{\mathring{\ll}} X$ et $Y'
\mathrel{\mathring{\ll}} Y$, alors $X' < Y'$.}
\end{quote}

En effet tout élément de $|X'| \subset |X|$ est $\leq$ tout élément de $|Y'|
\subset |Y|$ ; si $X' = Y'$, il est contenu dans $|X| \cap |Y|$, qui a au plus un
point, $\operatorname{ex} X = \operatorname{or} Y$. Dans ce cas $X'
\trianglelefteq X$ et $Y' \trianglelefteq Y$ ; avec $\mathrel{\mathring{\ll}}$,
cela force $X' = X$ et $Y' = Y$, donc $X = Y$, contre $X < Y$.

\textbf{Définition.} $X \mathrel{|\circ|} Y$ si $X < Y$ ou $Y < X$ : deux strates
sont disjointes quand elles sont comparables pour l'ordre « à gauche de ». Le
corollaire est l'axiome At~3, et At~4 est immédiat puisque $a < I_{a,b} < b$.
Les marges notent que At~1 et At~2 sont vérifiés.

\subsection*{Les préfigures sont les chaînes}

\begin{quote}
\textbf{Proposition} (page 6). \emph{Deux strates $X$, $Y$ de $\mathcal{A}(L)$
sont compatibles si et seulement si $X = Y$ ou $X \mathrel{|\circ|} Y$,
c'est-à-dire si $\{X, Y\}$ est un drapeau de $(\mathcal{A}(L), <)$.}
\end{quote}

La compatibilité est celle du chapitre VIII\footnote{La page note la relation
par un signe que la transcription rend par $\asymp$ (tracé barré) ; nous y
reconnaissons la relation de compatibilité du chapitre VIII (axiome At~M~5 du
dossier 156-8) : $X$ et $Y$ sont compatibles si toute sous-strate de $X$ qui
n'est pas sous-strate de $Y$ est disjointe de toute sous-strate de $Y$ qui n'est
pas sous-strate de $X$. La démonstration de la page est exactement celle de cette
définition. L'identification est la nôtre ; nous écrivons $\asymp$.} :
pour $X \neq Y$, $\widetilde{X} \cap \widetilde{Y} = \delta X \cap \delta Y$ a
au plus un élément. S'il est vide, la compatibilité dit que toutes les faces de
$X$ sont disjointes de toutes celles de $Y$, ce qui équivaut à $X
\mathrel{|\circ|} Y$. S'il a un élément $b$, ou bien l'une des strates est
incidente à l'autre, et les deux conditions sont vraies ; ou bien $X = I_{a,b}$
et $Y = I_{b,c}$ avec $a < b < c$, deux intervalles adjacents, et les deux
conditions sont encore vraies. Il en résulte :

\begin{quote}
\textbf{Corollaire.} \emph{Une partie finie de $\mathcal{A}(L)$ est une
préfigure si et seulement si elle est totalement ordonnée par $<$ ; les figures
sont les chaînes finies fermées pour $\trianglelefteq$.}
\end{quote}

Une figure, ici, c'est donc une suite $X_{1} < X_{2} < \cdots < X_{k}$ de points
et d'intervalles rangés de gauche à droite, avec les extrémités de chaque
intervalle : un sous-segment subdivisé, ou plusieurs morceaux successifs. On
écrit $\operatorname{\text{Préfig}} = \operatorname{Drap}^{*}(\mathcal{A}(L),
<)$, les chaînes finies non vides, et l'on travaille, comme la page 7 le
précise, dans le magasin des figures finies.

\pagerange{7}{15}
\section*{Intervalles infinis et préfigure complémentaire (pages 7 à 15)}

\subsection*{Pourquoi des intervalles infinis}

Pour étudier les supports et le passage au complémentaire, il veut associer à
toute préfigure une « préfigure complémentaire ». Si $L$ n'a pas de plus petit
ou de plus grand élément\footnote{La page écrit « Comme $L$ n'a pas de plus
petit \emph{et} de plus grand élément » ; la suite montre qu'il faut entendre
« n'a pas nécessairement ».}, ce qui est à gauche de la première strate d'une
figure n'est réuni par aucune strate. D'où des intervalles infinis $I_{<a}$,
$I_{>a}$. Même quand $L$ a des extrémités, ils rendent la construction plus
élégante : le complémentaire $\mathfrak{F}'$ d'une préfigure vérifie alors
$\delta \mathfrak{F}' \subset \delta \mathfrak{F}$, au lieu de $\delta
\mathfrak{F}' \subset \delta \mathfrak{F} \cup \{\omega_{0}, \omega_{1}\}$.

\subsection*{L'atelier élargi $\widehat{\mathcal{A}}(L)$}

On pose $\widehat{\mathcal{A}}(L) = \mathcal{A}(L) \sqcup \mathcal{A}_{1}^{-}
\sqcup \mathcal{A}_{1}^{+}$, avec
\[
\mathcal{A}_{1}^{-} = \{ I_{<a} \mid a \text{ non minimal dans } L\}, \qquad
\mathcal{A}_{1}^{+} = \{ I_{>a} \mid a \text{ non maximal dans } L\} .
\]
Ce sont des symboles, que l'on interprète comme la demi-droite ouverte à gauche
de $a$ : $\partial I_{<a} = \delta I_{<a} = \{a\}$, $|I_{<a}| = L_{\leq a}$,
$|I_{<a}|^{\circ} = L_{<a}$, et de même pour $I_{>a}$. Cela prolonge
$\trianglelefteq$ et $\lhd$, et l'on prolonge $\mathrel{\mathring{\ll}}$ par
\[
I_{<a} \mathrel{\mathring{\ll}} I_{<b} \iff a \leq b, \qquad X
\mathrel{\mathring{\ll}} I_{<a} \iff X < a \quad (X \in \mathcal{A}(L)),
\]
et symétriquement à droite\footnote{La page écrit « $X \leq a$ ». Avec la
définition générale qu'elle rappelle en marge --- $X \mathrel{\mathring{\ll}} Y$
si $|X| \subset |Y|$ et $X$ n'est pas incidente à $Y$ --- le cas $X = a$ est
exclu, puisque $a \lhd I_{<a}$. Il reste $X < a$.}. Deux remarques fixent
l'intention :
\begin{enumerate}
\item[a)] On n'a jamais $I_{<a} \mathrel{\mathring{\ll}} X$ pour $X \in
\mathcal{A}(L)$. Si $L$ a un plus petit élément $\omega_{0}$, on a
$I_{\omega_{0},a} \mathrel{\mathring{\ll}} I_{<a}$ et non l'inverse, bien que
les deux strates aient le même support $[\omega_{0}, a]$ : l'intervalle fini est
une « réalisation effective » de l'intervalle infini, qui ne lui est pas
identifié.
\item[b)] On n'a jamais $I_{<a} \mathrel{\mathring{\ll}} I_{>b}$, ni l'inverse.
\end{enumerate}

L'ordre se prolonge par : $X \leq Y$ si $\delta X \leq \delta Y$ (tout élément
de l'un est $\leq$ tout élément de l'autre), avec la restriction que rien n'est
strictement à gauche de $I_{<a}$ ni strictement à droite de $I_{>a}$. On a
toujours $X < Y \Rightarrow |X| < |Y|$, jamais $I_{>a} \leq I_{<b}$, et $I_{<a}
\leq I_{>b}$ si et seulement si $a \leq b$. La disjonction reste « être
comparables » : $X \mathrel{|\circ|} Y$ si $X < Y$ ou $Y < X$. Les axiomes At~1
à At~4 se vérifient comme pour $\mathcal{A}(L)$, avec le même corollaire.

\begin{quote}
\textbf{Lemme} (page 11). \emph{Soient $X \in \widehat{\mathcal{A}}(L)$ et $Y \in
\mathcal{A}(L)$. Pour que $X \mathrel{|\circ|} Y$, il faut et il suffit que $X'
\mathrel{|\circ|} Y$ pour toute strate $X' \in \mathcal{A}(L)$ telle que $X'
\mathrel{\mathring{\ll}} X$.}
\end{quote}

La nécessité est At~3\footnote{La page écrit « il faut », et démontre la
suffisance, seule non triviale. Elle énonce le lemme pour $X$ et $Y$ dans
$\widehat{\mathcal{A}}(L)$, mais suppose dans la démonstration $Y \in
\mathcal{A}(L)$ (« supposons d'abord ») ; le cas de deux intervalles infinis
n'est pas traité. Une première rédaction, biffée, visait une formule
$\operatorname{cosupp}^{\circ}(I_{<a}) = \{a\} \cup
\operatorname{Omb}^{\circ}(I_{>a})$.}. Pour la suffisance, prenons $X = I_{<a}$,
de sorte que $X \mathrel{|\circ|} Y$ signifie $a \leq Y$. Comme $a$ n'est pas
minimal, il existe $a' < a$, et $I_{a',a} \mathrel{\mathring{\ll}} I_{<a}$ ;
donc $Y \mathrel{|\circ|} I_{a',a}$, c'est-à-dire $Y \geq a$ ou $Y \leq a'$.
Dans le second cas $Y \mathrel{\mathring{\ll}} I_{<a}$, et en prenant $X' = Y$
on trouverait $Y \mathrel{|\circ|} Y$. Le corollaire de la page dit que les
supports de $\widehat{\mathcal{A}}(L)$ correspondent ainsi à ceux de
$\mathcal{A}(L)$, « ceux qui nous intéressent en fait »\footnote{L'énoncé du
corollaire est serré au bas de la page et lu avec doute. Ce qu'il faut pour
l'établir est le corollaire 5 de la page 87 : $\mathcal{A}(L)$ est « dense »
dans $\widehat{\mathcal{A}}(L)$ au sens de cette page, et la restriction $S
\mapsto S \cap \mathcal{A}(L)$ identifie alors $\Sigma_{\widehat{\mathcal{A}}(L)}$
à $\Sigma_{\mathcal{A}(L)}$. La densité demandée là est exactement la propriété
de ce lemme, une fois traité le cas de deux intervalles infinis.}.

\subsection*{La préfigure complémentaire}

On travaille désormais avec les préfigures finies de
$\widehat{\mathcal{A}}(L)$, c'est-à-dire ses chaînes finies. Pour une telle
préfigure $\mathfrak{F}$, $T = \delta \mathfrak{F}$ est un drapeau de $L$, non
vide si $\mathfrak{F}$ l'est. À un drapeau $T = \{t_{1} < \cdots < t_{n}\}$ on
associe ses \emph{structures ensemblistes}
\[
[I_{<t_{1}}],\ t_{1},\ I_{t_{1},t_{2}},\ t_{2},\ \ldots,\ I_{t_{n-1},t_{n}},\
t_{n},\ [I_{>t_{n}}],
\]
les termes entre crochets ne figurant que si $t_{1}$ n'est pas minimal, et $t_{n}$
pas maximal. Elles forment la plus grande préfigure $\widehat{\mathfrak{F}}_{T}$
de drapeau $T$, la subdivision complète de $L$ par les points de $T$. Une
préfigure de drapeau $T$ est une partie $\mathfrak{F}$ de
$\widehat{\mathfrak{F}}_{T}$ telle que tout $t \in T$ soit dans $\mathfrak{F}$
ou extrémité d'une strate de $\mathfrak{F}$.

\textbf{Définition} (page 13). La \emph{préfigure complémentaire} est
\[
\mathfrak{F}^{*} = \widehat{\mathfrak{F}}_{\delta\mathfrak{F}} \smallsetminus
\mathfrak{F}.
\]

Un sommet $s \in \mathfrak{F}$ est \emph{redondant} si toutes les strates de
$\widehat{\mathfrak{F}}_{T}$ qui lui sont adjacentes sont dans $\mathfrak{F}$ :
ses deux intervalles voisins, en général\footnote{La page définit
$\mathfrak{F}_{\mathrm{red}} = \{s \in \mathfrak{F}_{0} \mid
\operatorname{ord}(s, \mathfrak{F}) = 2\}$, le nombre d'intervalles de
$\mathfrak{F}$ issus de $s$. C'est la même chose tant que $s$ a deux voisins
dans $\widehat{\mathfrak{F}}_{T}$, c'est-à-dire sauf si $s$ est le plus petit ou
le plus grand élément de $L$ ; en ce cas il n'a qu'un voisin, et c'est la
définition que nous donnons qui rend juste la formule qui suit. Les pages 16 à
19 font cette distinction explicitement (sommets « semi-évanescents »).}. Alors
\[
\delta \mathfrak{F}^{*} = \delta \mathfrak{F} \smallsetminus
\mathfrak{F}_{\mathrm{red}},
\]
car un sommet redondant n'a plus, dans $\mathfrak{F}^{*}$, ni lui-même ni ses
voisins. La page renvoie, pour une classification des sommets en cinq types, à
des pages 110 et 111 de sa main qui ne sont pas dans ce dossier ; la seconde
présentation (pages 16 à 19) en refait une. $\mathfrak{F}^{*}$ n'a pas de sommet
redondant --- un sommet de $\mathfrak{F}^{*}$ dont tous les voisins seraient
dans $\mathfrak{F}^{*}$ ne serait pas dans $\delta\mathfrak{F}$ --- donc $\delta
\mathfrak{F}^{**} = \delta \mathfrak{F}^{*}$, et
\[
\mathfrak{F}^{**} = \{ X \text{ structure ensembliste de } \delta\mathfrak{F}
\smallsetminus \mathfrak{F}_{\mathrm{red}} \mid \exists\, Y \in \mathfrak{F},\
Y \mathrel{\mathring{\ll}} X \}
\]
s'obtient de $\mathfrak{F}$ « en effaçant les sommets redondants », c'est-à-dire
en composant les intervalles adjacents à ces sommets\footnote{La page écrit $Y
\ll X$. Avec $\ll$, qui contient l'incidence, un sommet isolé $a$ de
$\mathfrak{F}$ ferait entrer dans $\mathfrak{F}^{**}$ les intervalles dont il
est extrémité ; c'est $\mathrel{\mathring{\ll}}$ qu'il faut.}. Ses sommets sont
$\mathfrak{F}_{0} \smallsetminus \mathfrak{F}_{\mathrm{red}}$, et ses
intervalles sont les composés $I_{\operatorname{or} X_{1}, \operatorname{ex}
X_{k}}$ des suites maximales $X_{1} < x_{1} < X_{2} < \cdots < x_{k-1} < X_{k}$
contenues dans $\mathfrak{F}$, avec $x_{i} \lhd X_{i}, X_{i+1}$ ; elles
correspondent aux composantes connexes de $\mathfrak{F}$.

\subsection*{Le théorème}

\begin{quote}
\textbf{Théorème} (page 15). \emph{Pour toute préfigure non vide $\mathfrak{F}$
de $\widehat{\mathcal{A}}(L)$,}
\[
\operatorname{cosupp}^{\circ}_{\mathcal{A}(L)}(\mathfrak{F}) =
\operatorname{cosupp}^{\circ}\bigl(\operatorname{Omb}^{\circ}_{\mathcal{A}(L)}(\mathfrak{F})\bigr)
= \operatorname{Omb}^{\circ}_{\mathcal{A}(L)}(\mathfrak{F}^{*}) .
\]

\textbf{Corollaire 1.} \emph{Si $\mathfrak{F}^{*}$ n'est pas vide,}
\[
\operatorname{supp}^{\circ}_{\mathcal{A}(L)}(\mathfrak{F}) =
\operatorname{supp}^{\circ}\bigl(\operatorname{Omb}^{\circ}(\mathfrak{F})\bigr) =
\operatorname{Omb}^{\circ}(\mathfrak{F}^{**}) .
\]
\end{quote}

Ce qui est disjoint d'une préfigure, c'est ce qui est contenu dans les strates
ouvertes de sa complémentaire ; et le support d'une préfigure est l'ombre de la
préfigure qu'on obtient en effaçant ses sommets redondants\footnote{La page ne
démontre rien, et les hypothèses qu'elle avait d'abord écrites sur les
extrémités de $L$ sont biffées. Le corollaire 1 est énoncé « sous la condition
$\mathrm{H}_{\mathcal{L}}$ » --- $L$ a un plus petit élément ou n'a pas
d'élément minimal, et de même à droite ---, d'une lecture incertaine ; il suit
du théorème appliqué deux fois, à $\mathfrak{F}$ puis à $\mathfrak{F}^{*}$, ce
qui demande seulement $\mathfrak{F}^{*} \neq \emptyset$. Le corollaire 2
commence (« Soit $T \in \operatorname{Drap}(L)$. Considérons l'ens.\ des figures
telles que ») et s'interrompt ; la page 16 s'ouvre sur « C.~facile ».}. Voici
pourquoi le théorème est vrai. La première égalité est le lemme de la page 11 :
une strate de $\mathcal{A}(L)$ disjointe de l'ombre d'une strate infinie lui est
disjointe. Pour la seconde, écrivons $\mathfrak{F} = \{X_{1} < \cdots <
X_{k}\}$. Si $Z \in \mathcal{A}(L)$ est disjointe de chaque $X_{j}$, elle est
comparable à chacune, donc logée dans un « trou » : $X_{j} < Z < X_{j+1}$, ou
avant $X_{1}$, ou après $X_{k}$. Si $\operatorname{ex} X_{j} = t$ et
$\operatorname{or} X_{j+1} = t'$, on a $\delta Z \subset [t, t']$ ; si $t = t'$,
$Z = t$, qui n'est pas dans $\mathfrak{F}$ et est donc dans $\mathfrak{F}^{*}$ ;
si $t < t'$, ce sont deux points consécutifs de $T$, $I_{t,t'}$ n'est pas dans
$\mathfrak{F}$, et $Z$ est soit $t$ ou $t'$ (dans $\mathfrak{F}^{*}$ s'ils ne
sont pas dans $\mathfrak{F}$), soit contenue dans la strate ouverte $I_{t,t'}$.
Avant $X_{1}$, $Z$ est $\operatorname{or} X_{1}$ ou contenue dans
$I_{<\operatorname{or} X_{1}}$ ; de même après $X_{k}$. Réciproquement, toute
strate de $\mathfrak{F}^{*}$ est comparable à toute strate de $\mathfrak{F}$,
puisque $\widehat{\mathfrak{F}}_{T}$ est une chaîne, et At~3 transmet la
disjonction à son ombre.

Le nom actuel de ce que décrit cette section est celui d'\emph{ensemble
constructible} de la droite : une réunion finie de points et d'intervalles
ouverts\footnote{Les parties d'un ensemble totalement ordonné dense qui sont
réunions finies de points et d'intervalles sont exactement les parties
définissables en dimension 1 d'une structure o-minimale (L.~van den Dries, 1984 ;
A.~Pillay et C.~Steinhorn, 1986), notion contemporaine de ces pages, et que van
den Dries a présentée comme une réponse au programme de « topologie modérée » de
l'\emph{Esquisse d'un programme}. Le rapprochement est le nôtre ; les feuillets
ne citent personne.}. Le théorème dit que le complémentaire d'un constructible
est constructible, et donne la figure qui le réalise.

\pagerange{16}{20}
\section*{Autre présentation : le segment orienté combinatoire (pages 16 à 20)}

« Pour y comprendre quelque chose », il recommence avec un \emph{segment orienté
combinatoire} (ou « pseudo-segment ») : un ensemble ordonné $I$ muni d'un plus
petit élément $\omega_{0}$, l'\emph{origine}, et d'un plus grand $\omega_{1}$,
l'\emph{extrémité}, avec $\omega_{0} < \omega_{1}$ ; on note $\partial I =
\{\omega_{0}, \omega_{1}\}$. Quand $I = \{\omega_{0}, \omega_{1}\}$, le segment
est \emph{indivisible}, et son atelier est celui du simplexe orienté type
$\Delta^{1}$\footnote{Lecture incertaine (« magasin \ldots\ singulier
simpliciale ») ; le sens est que les strates sont les deux sommets et l'arête.}.
On n'introduit plus d'intervalles infinis : les strates sont les $a \in I$ et
les $I_{a,b}$, et la préfigure complémentaire se prend dans la subdivision
complète $\Phi_{T}$ d'un drapeau $T$ \emph{contenant $\partial I$}, de sorte que
le complément ne sort jamais du segment.

\subsection*{Les positions d'un sommet}

Soit $\mathfrak{F}$ une préfigure, $T = \delta\mathfrak{F} \cup \partial I$, et
pour $s \in T$ notons $V(s)$ les intervalles de $\Phi_{T}$ adjacents à $s$ :
deux si $s$ est intérieur, un seul si $s \in \partial I$. Les pages 17 à 19
distinguent\footnote{Les pages essaient plusieurs noms à la suite, et le terme
« redondant » y désigne tour à tour l'ensemble des sommets qui disparaissent
par dualité et la partie de cet ensemble qui est dans $\partial I$. Nous
retenons le vocabulaire final de la page 19 (« évanescents », « semi-évanescents »)
et donnons les définitions qui rendent justes les formules des pages 18 et 19.} :
\begin{itemize}
\item $s$ \emph{isolé} : $s \in \mathfrak{F}$, aucun intervalle de $V(s)$ dans
$\mathfrak{F}$ ;
\item $s$ \emph{lacunaire} : $s \notin \mathfrak{F}$, tous les intervalles de
$V(s)$ dans $\mathfrak{F}$ --- un trou ponctuel ;
\item $s$ \emph{évanescent} : $s \notin \partial I$, $s \in \mathfrak{F}$ avec
ses deux voisins --- c'est le sommet redondant de la page 13 ;
\item $s$ \emph{semi-évanescent} : $s \in \partial I$, $s \in \mathfrak{F}$ avec
son unique voisin ; et $s$ est semi-évanescent \emph{pour $\mathfrak{F}^{*}$}
si $s \in \partial I$ et que ni $s$ ni son voisin ne sont dans $\mathfrak{F}$ ;
\item $s$ \emph{sommet-bord} : $s \notin \partial I$, un seul de ses deux
voisins dans $\mathfrak{F}$ ; \emph{propre} si $s \in \mathfrak{F}$, \emph{impropre}
sinon. On note $\partial^{*}\mathfrak{F}$ leur ensemble.
\end{itemize}
On obtient une décomposition disjointe
\[
\widehat{\delta}\mathfrak{F} = \delta\mathfrak{F} \cup \partial I =
\mathfrak{F}_{\mathrm{is}} \sqcup \mathfrak{F}_{\mathrm{lac}} \sqcup
\mathfrak{F}_{\text{év}} \sqcup \mathfrak{F}_{\text{sév}}
\sqcup \mathfrak{F}^{*}_{\text{sév}} \sqcup \partial^{*}\mathfrak{F},
\]
et le passage à la complémentaire agit ainsi :
\[
\mathfrak{F}^{*}_{\mathrm{is}} = \mathfrak{F}_{\mathrm{lac}}, \quad
\mathfrak{F}^{*}_{\mathrm{lac}} = \mathfrak{F}_{\mathrm{is}}, \quad
\partial^{*}\mathfrak{F}^{*} = \partial^{*}\mathfrak{F}, \quad
\mathfrak{F}^{*}_{\text{év}} = \emptyset ,
\]
les semi-évanescents de $\mathfrak{F}$ et de $\mathfrak{F}^{*}$ s'échangent, et
les sommets-bords propres et impropres aussi. Les évanescents disparaissent par
dualité et ne reviennent pas ; les semi-évanescents disparaissent de
$\delta\mathfrak{F}^{*}$ mais « réapparaissent par bidualité », puisque
$\partial I$ fait toujours partie du drapeau. C'est pourquoi le mot
« redondant » ne leur convient pas, dit la page 18. Pour un point $s$ de
$\partial I$, il y a exactement quatre positions, deux à deux échangées.

\subsection*{Parties constructibles}

On retrouve la formule encadrée
\[
\operatorname{cosupp}^{\circ}(\mathfrak{F}) =
\operatorname{cosupp}^{\circ}\bigl(\operatorname{Omb}^{\circ}(\mathfrak{F})\bigr) =
\operatorname{Omb}^{\circ}(\mathfrak{F}^{*}) .
\]
Appelons \emph{constructibles}\footnote{La transcription lit « contractibles »
(pages 7, 19 et 20). Le nom de l'ensemble, lu « $\operatorname{Cons}$ » avec
doute, et l'usage des pages 21 à 26 (« supports constructibles ») font préférer
« constructibles » ; l'ombre d'une préfigure n'est pas connexe en général, et
« contractible » ne lui convient pas. Les deux lectures sont possibles sur la
page ; c'est un choix de lecture.} les parties de $\mathcal{A}(I)$ de la forme
$\operatorname{Omb}^{\circ}(\mathfrak{F})$. L'application $\mathfrak{F} \mapsto
\operatorname{Omb}^{\circ}(\mathfrak{F})$ des préfigures vers les parties
constructibles est bijective, et l'on retrouve $\mathfrak{F}$ à partir de
$\Sigma = \operatorname{Omb}^{\circ}(\mathfrak{F})$ :
\begin{itemize}
\item $\mathfrak{F}_{1}$ est l'ensemble des intervalles de $\Sigma$ maximaux
pour $\mathrel{\mathring{\ll}}$ ;
\item $\mathfrak{F}_{\mathrm{is}}$ est l'ensemble des points de $\Sigma$
maximaux pour $\ll$\footnote{La page écrit « maximaux pour
$\mathrel{\mathring{\ll}}$ ». Mais tout sommet de $\mathfrak{F}$ est maximal pour
$\mathrel{\mathring{\ll}}$ dans $\Sigma$, puisqu'il n'est intérieur à aucune
strate de $\mathfrak{F}$ ; seuls les sommets isolés sont maximaux pour $\ll$, qui
voit aussi l'incidence. C'est ce qui rend la troisième ligne nécessaire.} ;
\item les autres sommets de $\mathfrak{F}$ sont les extrémités des intervalles de
$\mathfrak{F}_{1}$ qui sont dans $\Sigma$.
\end{itemize}
Le cosupport d'une partie constructible est donc constructible, et le passage au
cosupport correspond, du côté des préfigures, au passage à la complémentaire.

Pour un drapeau $T \supset \partial I$, $\omega_{0} = t_{0} < t_{1} < \cdots <
t_{n} = \omega_{1}$, la subdivision complète est
\[
\Phi_{T} = \{ t_{0},\ I_{t_{0},t_{1}},\ t_{1},\ \ldots,\ I_{t_{n-1},t_{n}},\
t_{n} \},
\]
et une \emph{préfigure subordonnée à $T$} est une partie quelconque $\Phi$ de
$\Phi_{T}$ ; on note $\complement\Phi = \Phi_{T} \smallsetminus \Phi$ son
complémentaire dans $\Phi_{T}$\footnote{La fin de la page 20, écrite très vite,
est d'une lecture très incertaine ; elle annonce un « principe de dualité » pour
ces parties, que la page 21 formule.}.

\pagerange{21}{27}
\section*{Supports subordonnés à un drapeau ; segments ouverts (pages 21 à 27)}

\subsection*{L'algèbre de Boole d'un drapeau}

Pour $\Phi \subset \Phi_{T}$, la complémentaire $\complement\Phi$ dans
$\Phi_{T}$ ne diffère de $\Phi^{*}$ que par des sommets évanescents, et
$(\complement\Phi)^{**} = \Phi^{*}$. Avec le théorème et son corollaire, il
vient
\[
\boxed{\ \operatorname{cosupp}^{\circ}\Phi = \operatorname{Omb}^{\circ}(\Phi^{*})
= \operatorname{supp}^{\circ}\operatorname{Omb}^{\circ}(\complement\Phi) =
\operatorname{supp}^{\circ}(\complement\Phi)\ }
\]
(page 21). Posons $S_{\Phi} = \operatorname{supp}^{\circ}(\Phi)$. L'application
\[
\mathfrak{P}(\Phi_{T}) \longrightarrow \Sigma_{\mathrm{cons}}(\mathcal{A}(I)),
\qquad \Phi \longmapsto S_{\Phi},
\]
vers l'ensemble des supports constructibles, envoie les réunions sur les Sup (par
définition) ; la formule encadrée dit qu'elle commute à $\complement$ ; elle
commute donc aux intersections\footnote{La page ajoute que cela « utilise
l'axiome des supports Ax supp 1 », d'une lecture incertaine ; le raisonnement
ci-dessus n'en a pas besoin.}. Elle est de plus injective, puisque $S_{\Phi} \cap
\Phi_{T} = \Phi$ : un sommet $t_{i}$ est dans $\operatorname{Omb}^{\circ}(\Phi^{**})$
si et seulement s'il est dans $\Phi$, et un intervalle $I_{t_{i},t_{i+1}}$ de
même\footnote{L'injectivité n'est pas énoncée, mais la page compte
$\operatorname{card}\mathfrak{P}(\Phi_{T}) = 2^{2n+1}$ comme « nombre de supports
subordonnés à $T$ », ce qui la suppose ; la justification est la nôtre.}. Si $T$
a $n+1$ points, $\Phi_{T}$ a $2n+1$ éléments, et \emph{les supports subordonnés à
$T$ forment une algèbre de Boole à $2^{2n+1}$ éléments, isomorphe à
$\mathfrak{P}(\Phi_{T})$}.

Deux préfigures $\Phi$, $\Psi$ sont \emph{commensurables} si $\delta\Phi \cup
\delta\Psi$ est un drapeau. Leurs supports sont alors subordonnés au même
drapeau $T = \delta\Phi \cup \delta\Psi \cup \partial I$, quitte à subdiviser les
intervalles, ce qui ne change pas les supports. Conclusion : toutes les
opérations usuelles sur une famille de supports constructibles mutuellement
commensurables se ramènent à $\cup$, $\cap$, $\complement$ dans un
$\mathfrak{P}(\Phi_{T})$ convenable. Cette restriction --- on ne calcule
qu'entre supports commensurables --- est celle que la critique des pages 123 à
126 opposera à la théorie générale.

\subsection*{Segments sans extrémités données}

Si l'on ne se donne pas $\omega_{0}$ et $\omega_{1}$, on voit le segment comme
obtenu du cas précédent en ôtant une extrémité, ou les deux, devenue
« virtuelle » :
\begin{itemize}
\item $\mathcal{A}(J, \omega_{1}) = \mathcal{A}(I, \partial I) \smallsetminus
\{\omega_{0}\}$, avec $J = I \smallsetminus \{\omega_{0}\}$ : les strates
$I_{\omega_{0},a}$ deviennent les $J_{<a}$, et l'on a les strates $a$,
$J_{<a}$, $J_{a,b}$ ;
\item dualement $\mathcal{A}(J, \omega_{0})$, avec les $J_{>a}$ ;
\item $\mathcal{A}(J) = \mathcal{A}(I, \partial I) \smallsetminus \partial I$,
avec $J = I \smallsetminus \partial I$ : les $J_{<a}$, les $J_{>a}$, les
$J_{a,b}$, et une strate « unité » $[J]$, trace de $I_{\omega_{0},\omega_{1}}$,
indispensable pour exprimer la subdivision complète d'un drapeau vide.
\end{itemize}
Ces strates se décrivent en termes de $J$ seul, sans référence à $I$. Deux
exemples. Si $I = \{\omega_{0}, \omega_{1}\}$, $\mathcal{A}(J, \omega_{1}) =
\{\omega_{1}, J_{<\omega_{1}}\}$ a deux strates, une maquette, et
$|J_{<\omega_{1}}|^{\circ}$ est vide bien que la strate existe. Et
$\mathcal{A}(\emptyset) = \{[\emptyset]\}$ n'a qu'une strate : « l'intervalle
ouvert indivisible est représenté par l'ensemble ordonné vide ! » --- cas qu'il
ne faut nullement exclure.

La formulation commune est un triplet $(I, \Omega_{0}, \Omega_{1})$ : un ensemble
ordonné et deux parties, $\Omega_{0}$ formée de plus petits éléments,
$\Omega_{1}$ de plus grands, avec $\Omega_{0} \cap \Omega_{1} = \emptyset$ ;
chacune a au plus un élément, d'où quatre cas. On peut le dire ainsi :
$\mathcal{A}(I, \Omega_{0}, \Omega_{1})$ est l'atelier du segment obtenu en
adjoignant à $I$ une origine virtuelle si $\Omega_{0} = \emptyset$ et une
extrémité virtuelle si $\Omega_{1} = \emptyset$, privé des points
adjoints\footnote{Cette formulation unifiée est la nôtre ; la page dit qu'on
« doit voir » chaque cas comme se déduisant de $(I, \partial I)$. On notera la
différence avec l'atelier $\widehat{\mathcal{A}}(L)$ des pages 8 à 10 : là,
$I_{<a}$ n'existe que pour $a$ non minimal, coexiste avec $I_{\omega_{0},a}$
quand $L$ a un plus petit élément, et il n'y a pas de strate unité. La seconde
présentation est la plus économe.}. Pour décrire un atelier et y faire la théorie
des préfigures, des parties constructibles et des supports, il n'y a pas d'autre
hypothèse à faire ; pour les préfigures subordonnées à un drapeau, on a intérêt
à prendre un drapeau contenant $\partial I = \Omega_{0} \cup \Omega_{1}$.

\subsection*{L'atelier induit sur une partie constructible}

L'atelier induit sur l'ombre d'une préfigure $\Phi$ est la somme (voir la page
28) des ateliers induits par ses composantes connexes, et pour $\Phi$ connexe il
est de la forme $\mathcal{A}(J', \Omega'_{0}, \Omega'_{1})$, où
\begin{itemize}
\item $J' = \operatorname{Omb}^{\circ}(\Phi) \cap L$ est la somme lexicographique
des intérieurs $|X|^{\circ}$, $X \in \Phi$, dans l'ordre de $\Phi$ ;
\item $\Omega'_{0}$ est vide si la première strate de $\Phi$ n'a pas d'origine
(c'est un $I_{<a}$) ou si son origine n'est pas dans $\Phi$, et
$\{\operatorname{or}\Phi\}$ sinon ; $\Omega'_{1}$ dualement.
\end{itemize}
Par exemple l'ombre de $\{I_{a,b}\}$ est l'atelier du segment ouvert $]a,b[$,
$I_{a,c}$ y devenant $J'_{<c}$ et $I_{a,b}$ la strate unité ; l'ombre de $\{a,
I_{a,b}\}$ est celui de $[a,b[$\footnote{La description de $J'$ par trois
conditions, au bas de la page 26, est serrée et en partie illisible, de même que
la mise en garde qui la suit (« Attention cette description \ldots\ si $\Phi$
\ldots ») ; la page 27 la remplace par la somme lexicographique. Il note qu'il
faut définir l'origine d'une préfigure, quand sa première strate n'est pas un
$I_{<a}$.}.

\subsection*{Divisibilité}

L'atelier est \emph{divisible} si, pour $a < b$, il existe $c$ avec $a < c < b$ ;
pour un segment ouvert d'un côté, il faut en outre que tout élément ait un
voisin strictement plus près du côté ouvert\footnote{La page dit : « dans le cas
de $(I, \omega_{0})$ resp.\ $(I, \omega_{1})$, il faut ajouter que $\forall a$
existe $b$ t.q.\ $b < a$ resp.\ $b > a$ ». Pour $(I, \omega_{0})$, ouvert à
droite, c'est $b > a$ qu'il faut (sinon $a = \omega_{0}$ contredit la
condition) : le « resp. » est inversé.} ; pour un segment ouvert des deux
côtés, que $I \neq \emptyset$ --- bien que $\mathcal{A}(\emptyset)$, réduit à
une strate, soit divisible, « ce qui gêne ». Si $I$ est divisible, les strates
« compactes » $a$ et $I_{a,b}$ sont « très denses » dans l'atelier.

\pagerange{28}{40}
\section*{Quatre digressions (pages 28 à 40)}

\subsection*{1. Somme d'ateliers (pages 28 et 29)}

La somme d'une famille $(\mathcal{A}_{i})$ est la réunion disjointe, avec les
relations $\trianglelefteq$ et $\mathrel{\mathring{\ll}}$ sommes, et $X
\mathrel{|\circ|} Y$ si $X$ et $Y$ sont dans des termes différents ou disjoints
dans le même terme ; les axiomes se vérifient tautologiquement. Les magasins
$\mathcal{F}_{i}$ donnent le magasin produit $\mathcal{F} = \{A \mid A \cap
\mathcal{A}_{i} \in \mathcal{F}_{i}\ \forall i\}$\footnote{La page écrit
« $\forall i \in \mathcal{M}_{i}$ » ; il faut lire $\forall i$. Pour une famille
infinie, ce produit contient des figures infinies.}. Ce n'est pas une somme au
sens des catégories pour les morphismes considérés jusqu'ici, mais elle
correspond à la somme des espaces. Inversement les composantes connexes d'un
atelier pour $\ll$ en donnent une décomposition en somme, sous l'axiome qu'il
nomme « Mag connexe » : deux strates de composantes différentes sont disjointes.

\subsection*{2. Sous-ateliers (pages 30 à 32)}

Une partie $\mathcal{A}'$, avec les relations induites, est un atelier dès
qu'elle satisfait la condition

\textbf{Sous-at.} \emph{Si $X \mathrel{\mathring{\ll}} Y' \trianglelefteq Y$ avec
$X, Y \in \mathcal{A}'$, alors $Y' \in \mathcal{A}'$}

--- c'est la forme que la marge propose, plus forte que celle du
texte\footnote{Le texte demande seulement que, si $X \mathrel{\mathring{\ll}} Y$
et $X' \trianglelefteq X$ dans $\mathcal{A}'$, la strate $Y' \trianglelefteq Y$
telle que $X' \mathrel{\mathring{\ll}} Y'$ soit dans $\mathcal{A}'$ ; c'est
automatique si $\mathcal{A}'$ est fermée pour $\trianglelefteq$. La marge, en
oblique et en partie illisible, donne la forme ci-dessus, qu'il reprend
littéralement au chapitre IX bis, page 129.}.

\emph{Exemple 1.} L'ombre $\operatorname{Omb}^{\circ}(\Phi)$ d'une préfigure est
un sous-atelier --- fermée pour $\mathrel{\mathring{\ll}}$, mais pas pour
$\trianglelefteq$ en général. La démonstration (page 31) repose sur l'unicité :
une strate est contenue dans au plus une strate ouverte d'une figure, puisque
deux strates distinctes d'une figure sont disjointes et que At~3 interdirait
alors $X' \mathrel{|\circ|} X'$.

\emph{Exemple 2.} Le cosupport d'une préfigure n'est pas un sous-atelier en
général. Le contre-exemple de la page 31 est dessiné dans un triangle $abc$,
$d$ étant un point de l'arête ouverte $]b,c[$ : $Y = (a,b,c)$, $X = (a,b,d)$,
$X' = (b,d)$, $Y' = (b,c)$, $\Phi = \{d\}$. Le point $d$ est disjoint des strates
ouvertes $X$, $Y$ et $X'$, dont il est un sommet ou un point du bord, mais pas
de l'arête ouverte $Y'$ ; or $X \mathrel{\mathring{\ll}} Y$, $X' \trianglelefteq
X$, $X' \mathrel{\mathring{\ll}} Y' \trianglelefteq Y$. « Donc, attention à la
définition du magasin visible d'une figure. »

Il propose alors de se borner à
\[
\mathcal{A}' = \{ X \in \mathcal{A} \mid X \mathrel{|\circ|} \Phi \text{ et }
\Phi \text{ n'est pas une face de } X \},
\]
le « magasin visible hors de $\Phi$ », et affirme que c'est un
sous-atelier\footnote{C'est faux au-delà de la dimension 1. Gardons le triangle
et prenons un point $e$ de $]b,d[$ : $X = (a,b,e)$, $X' = (b,e)$, $Y = (a,b,c)$
sont disjoints de $d$ sans l'avoir pour sommet, $X \mathrel{\mathring{\ll}} Y$,
$X' \trianglelefteq X$, $X' \mathrel{\mathring{\ll}} Y' = (b,c) \trianglelefteq
Y$, et $Y'$ n'est pas disjoint de $d$. La démonstration de la page ne traite que
le cas $Y' \in \Phi$. Dans les ateliers de dimension 1 de ce dossier,
l'énoncé est vrai pour une raison simple : $Y'$ y est $Y$ ou $X'$. La version
juste est au chapitre IX bis (page 130), où la condition « $\Phi$ n'est pas une
face de $X$ » est remplacée par « $X$ est compatible avec $\Phi$ ». La page écrit
« $X, Y, X' \in \mathcal{M} \Rightarrow Y' \in \mathcal{M}$ » pour
$\mathcal{A}'$.}.

\subsection*{3. Ateliers compactologiques (pages 33 à 36)}

On se donne en plus une partie $\mathcal{A}_{c}$ de strates dites
\emph{compactes}, avec l'axiome
\begin{enumerate}
\item[a)] si $X$ est compacte, $\widetilde{X}$ est fini ;
\item[b)] $\mathcal{A}_{c}$ est fermée vers le bas pour $\ll$ ;
\item[c)] pour toute strate $X$, la famille $\widetilde{X}$ est localement finie
(page 34).
\end{enumerate}
Une partie $\Phi$ est \emph{localement finie} si chaque strate compacte ne
rencontre (adhérences comprises) qu'un nombre fini d'éléments de
$\Phi$\footnote{La page note $Y \mathbin{\overline{\|}} X$ la négation de $Y
\mathbin{\|} X$, glosé à la page 33 « $\widetilde{Y} \mathrel{|\circ|}
\widetilde{X}$ » : les adhérences sont disjointes.}. Les figures localement
finies forment un magasin $\mathcal{F}_{\mathrm{lf}}$ ; un magasin est
compatible avec la compactologie (axiome « Mag at cpct ») si et seulement s'il
est contenu dans $\mathcal{F}_{\mathrm{lf}}$, qui est donc le plus grand
possible. Inversement, un magasin donné détermine une plus grande compactologie
compatible. Il annonce un futur axiome de type Borel--Lebesgue, et note que
$\mathcal{A}_{c}$ devrait être « riche ». Il ne lui paraît pas raisonnable de
définir $\mathcal{A}_{c}$ à partir de l'atelier seul : pour un ensemble ordonné
sans extrémité, les strates naturellement compactes sont les seules $a$ et
$I_{a,b}$ --- les intervalles dont l'adhérence est un segment ---, alors que le
critère intrinsèque les déclarerait toutes compactes. Conclusion : la bonne
donnée est un \emph{atelier compactologique} $(\mathcal{A}, \trianglelefteq,
\mathrel{\mathring{\ll}}, \mathrel{|\circ|}, \mathcal{A}_{c})$, plus simple et
plus souple qu'un magasin particulier ; en pratique on n'utilisera que les
magasins des figures finies, des figures localement finies, ou de toutes les
figures. Les figures compactes forment un magasin attenant à $\mathcal{A}_{c}$,
non à $\mathcal{A}$ sauf si $\mathcal{A}_{c} = \mathcal{A}$\footnote{Une note
marginale (page 33) ajoute qu'un atelier est « à strates compactes » si
$\mathcal{A}_{c} = \mathcal{A}$, et \emph{compact} s'il existe une figure finie
de support total.}.

\subsection*{4. Terminologie (pages 37 à 40)}

« De plus en plus, l'accent principal dans le travail de fondements se porte sur
$\mathcal{M}$, plutôt que sur $\mathcal{F}$. » D'où un échange de noms :
l'\emph{atelier} est l'ensemble des strates, qui « travaille » avec elles pour
constituer des figures ; un ensemble de figures est un \emph{magasin attenant},
plusieurs magasins pouvant être attenants au même atelier. Dans un atelier
compactologique, $\mathcal{A}_{c}$ serait le \emph{cœur}. Les éléments de
$\mathcal{A}$ sont les \emph{strates}, $\widetilde{X}$ est la \emph{multistrate}
associée à $X$ ; pour une figure $F$, les strates de $F$ sont les éléments de
$\widetilde{F}$. Dans l'atelier ensembliste d'un ensemble $L$, une strate est un
ensemble de parties de $L$ et une multistrate un ensemble d'ensembles de
parties. Mais $\mathcal{A}$ étant désormais l'ensemble de base, les figures en
sont des parties, et il n'y a plus lieu de parler de leur « déploiement » ; on
verra plus tard s'il faut rendre $\mathcal{F}$ autonome. « Dès maintenant, je
dirai “atelier” là où je disais “magasin”, et inversement. Les axiomes Mag~1 à
Mag~4 deviennent désormais At~1 à At~4. » C'est la convention de ce document
depuis le début.

\pagerange{40}{60}
\section*{Supports, ateliers spatiaux, spatios (pages 40 à 60)}

\subsection*{Le treillis des supports}

Soit un atelier $(\mathcal{A}, \trianglelefteq, \mathrel{\mathring{\ll}},
\mathrel{|\circ|})$. La relation symétrique et antiréflexive
$\mathrel{|\circ|}$ suffit à définir les \emph{supports} (« contours ») : les
parties de la forme $\operatorname{cosupp}^{\circ}(B)$, ensemble des strates
disjointes de toutes celles de $B$. Elles forment une partie
$\Sigma_{\mathcal{A}}$ de $\mathfrak{P}(\mathcal{A})$ stable par intersections
quelconques, de plus petit élément $\emptyset$ et de plus grand $\mathcal{A}$ ;
$\complement = \operatorname{cosupp}^{\circ}$ y est une involution qui renverse
l'ordre, d'où l'existence des Sup, et $\complement$ échange Inf et Sup. On
étend la disjonction aux supports, $S \mathrel{|\circ|} S'$ si tout élément de
l'un est disjoint de tout élément de l'autre, et
\[
S \mathrel{|\circ|} S' \iff S \subset \complement S' \iff S' \subset \complement
S, \qquad S \mathrel{|\circ|} S' \Longrightarrow S \cap S' = \emptyset,
\]
la réciproque étant fausse en général. Les supports sont fermés pour
$\mathrel{\mathring{\ll}}$, non pour $\trianglelefteq$. C'est la construction
classique des ensembles fermés d'une \emph{polarité}, c'est-à-dire d'une
connexion de Galois entre $\mathfrak{P}(\mathcal{A})$ et lui-même\footnote{G.~Birkhoff,
\emph{Lattice Theory} (1940), pour les polarités ; O.~Ore (1944) pour les
connexions de Galois. Le rapprochement est le nôtre.}.

Parmi les supports, les \emph{supports constructibles}
$\Sigma\mathrm{cons}_{\mathcal{A}}$ sont définis à l'aide des figures et de
$\trianglelefteq$ ; ils ne sont stables ni par Sup, ni par Inf, ni par
$\complement$, et un « axiome des supports », sur lequel il reviendra (pages 114
à 121), en tient lieu. Pour une strate $X$, son support $X^{\circ} =
\operatorname{supp}^{\circ}\{X\}$ est le plus petit support contenant $X$ : c'est
une explicitation de la « partie occupée » par $X$.

\subsection*{Déploiement et ateliers spatiaux}

Le \emph{déploiement} (spatial) d'une strate est
\[
\operatorname{\text{Dép}}(X) = \{ X'^{\circ} \mid X' \in \widetilde{X} \}
\subset \Sigma\mathrm{cons}_{\mathcal{A}} .
\]
Ses éléments sont deux à deux disjoints, puisque deux sous-strates distinctes de
$X$ le sont (At~4), et non nuls ; l'application $\widetilde{X} \to
\operatorname{\text{Dép}}(X)$ est donc bijective, et
$\operatorname{\text{Dép}}(X)$ hérite de $\widetilde{X}$ un ordre, de plus grand
élément $X^{\circ}$\footnote{Ce plus grand élément l'est pour l'ordre transporté
de $\widetilde{X}$, non pour l'inclusion dans $\Sigma$ : les supports des faces
de $X$ sont disjoints de $X^{\circ}$, les strates étant ouvertes.}. C'est une
« figure élémentaire ordonnée » du treillis $\Sigma$.

Si $X \ll Y$, At~2 donne une application croissante $\varphi : \widetilde{X} \to
\widetilde{Y}$ avec $X' \mathrel{\mathring{\ll}} \varphi(X')$ ; elle se traduit
par une application croissante $\psi : \operatorname{\text{Dép}}(X) \to
\operatorname{\text{Dép}}(Y)$ avec $S \subset \psi(S)$, et $\psi$ est déterminée
par cette propriété, car un support non nul ne peut être contenu dans deux
supports disjoints. D'où une condition nécessaire pour que $X \ll Y$ :

\textbf{(R)} \emph{tout élément de $\operatorname{\text{Dép}}(X)$ est contenu
dans un élément de $\operatorname{\text{Dép}}(Y)$, et l'application $\psi$ ainsi
définie est croissante}

--- et l'on a de plus $X \mathrel{\mathring{\ll}} Y$ si et seulement si
$X^{\circ} \subset Y^{\circ}$.

\begin{quote}
\textbf{Définition} (pages 45 et 46). Un atelier est \emph{spatial} si la
condition (R) est aussi suffisante pour que $X \ll Y$.
\end{quote}

Le mot veut suggérer que la relation primitive $\ll$ se décrit en termes de
$\trianglelefteq$ et de $\mathrel{|\circ|}$, par référence à la notion
d'« espace occupé » par une figure. D'où la question : quand une structure
$(\mathcal{A}, \trianglelefteq, \mathrel{|\circ|})$ provient-elle d'un atelier
spatial ? On construit $\Sigma$, $\complement$ et $\operatorname{\text{Dép}} :
\mathcal{A} \to \operatorname{\text{Figél ord}}(\Sigma)$, à valeurs dans les
parties ordonnées de $\Sigma$ à plus grand élément, formées d'éléments non nuls
deux à deux disjoints ; on définit $X \mathrel{\mathring{\ll}} Y$ par la condition
(R) avec $\psi$ envoyant plus grand élément sur plus grand élément. C'est un
préordre ; il faut qu'il soit un ordre, ce qui revient à l'injectivité de
$\operatorname{\text{Dép}}$ ; les axiomes At~1 à At~3 se vérifient alors, At~4
devant être imposé d'avance. D'où :

\begin{quote}
\textbf{Proposition} (page 49). \emph{Un atelier spatial est la même chose qu'une
structure $(\mathcal{A}, \trianglelefteq, \mathrel{|\circ|})$ satisfaisant}
\end{quote}
\begin{enumerate}
\item[Atspat 1.] $\trianglelefteq$ \emph{est un ordre, $\mathrel{|\circ|}$ est
symétrique et antiréflexive ;}
\item[Atspat 2.] \emph{si $X \trianglelefteq Y$ et $X \neq Y$, alors $X
\mathrel{|\circ|} Y$}\footnote{La page écrit « Si $X \trianglelefteq Y$, on a $X
\mathrel{|\circ|} Y$ », ce qui contredirait l'antiréflexivité pour $X = Y$ ; il
faut $X \neq Y$, ce qui est At~4.} ;
\item[Atspat 3.] \emph{l'application $X \mapsto \operatorname{\text{Dép}}(X)$
(avec l'ordre transporté de $\widetilde{X}$) est injective, où $Y^{\circ} = \{ Z
\mid \forall Z',\ Z' \mathrel{|\circ|} Y \Rightarrow Z' \mathrel{|\circ|} Z
\}$.}
\end{enumerate}

\subsection*{Les spatios}

\begin{quote}
\textbf{Définition} (pages 49 et 50). Un \emph{spatios} est un triplet $(\Sigma,
\leq, \complement)$ tel que : Spat~1, $\leq$ est un ordre admettant des Inf (donc
des Sup) quelconques, de plus petit élément $0_{\Sigma}$ et de plus grand
$1_{\Sigma}$ ; Spat~2, $\complement$ est une involution qui renverse l'ordre ;
Spat~3, $x \wedge \complement x = 0_{\Sigma}$ pour tout $x$ (ou, ce qui revient
au même, $x \vee \complement x = 1_{\Sigma}$).
\end{quote}

Le nom actuel est \emph{treillis orthocomplémenté complet}, ou ortho-treillis
complet\footnote{L'ortho-complémentation est l'involution qui renverse l'ordre
avec $x \wedge x^{\perp} = 0$ ; la notion vient de G.~Birkhoff et J.~von Neumann,
« The logic of quantum mechanics » (1936). Les feuillets n'en parlent pas ; le
nom est le nôtre.}. On y pose $x \mathrel{|\circ|} y$ si $x \leq \complement y$ :
c'est symétrique, entraîne $x \wedge y = 0$, et $x \mathrel{|\circ|} x$ seulement
pour $x = 0$.

\emph{Exemple} (pages 50 et 51). Pour un ensemble $\mathcal{A}$ muni d'une
relation symétrique $\mathrel{|\circ|}$, l'ensemble de ses supports satisfait
Spat~1 et Spat~2, avec pour plus petit élément l'ensemble $\mathcal{A}_{00}$ des
éléments disjoints de tout ; il satisfait Spat~3 si et seulement si tout élément
disjoint de lui-même est disjoint de tout --- « la condition que j'avais mise dès
le début ». Appliquée à un spatios $\Sigma$ lui-même, la construction donne
$\operatorname{cosupp}(A) = \{ Z \mid Z \leq \complement\operatorname{Sup}A\}$ et
$\operatorname{supp}(A) = \{ Z \mid Z \leq \operatorname{Sup}A\}$ : les supports
de $\Sigma$ sont ses parties principales $\{Z \leq \xi\}$, et $\Sigma$ s'identifie
à l'ensemble de ses propres supports. C'est la première forme de la proposition
1 de la page 79.

\subsection*{L'atelier universel d'un spatios}

Une \emph{figure} d'un spatios est un couple $(\mathfrak{F}, \rho)$ d'une partie
$\mathfrak{F} \subset \Sigma$ et d'un ordre $\rho$ sur elle, dont les éléments
sont non nuls et deux à deux disjoints ; elle est \emph{élémentaire} si elle a un
plus grand élément. Sur l'ensemble $\operatorname{Atspat}(\Sigma)$ des figures
élémentaires, on met : $\trianglelefteq$ induit par l'atelier ensembliste de
$\Sigma$ ; $(\mathfrak{F}, \rho) \mathrel{\mathring{\ll}} (\mathfrak{F}', \rho')$
s'il existe $\psi : \mathfrak{F} \to \mathfrak{F}'$ croissante avec $X \leq
\psi(X)$, envoyant plus grand élément sur plus grand élément ; et
$(\mathfrak{F}, \rho) \mathrel{|\circ|} (\mathfrak{F}', \rho')$ si leurs plus
grands éléments sont disjoints dans $\Sigma$. C'est un atelier spatial (« At~1
\ldots\ At~4 clair »), son spatios des supports s'identifie à $\Sigma$ par
$(\mathfrak{F}, \rho) \mapsto \max \mathfrak{F}$, et l'application canonique
$\operatorname{Atspat}(\Sigma) \to \operatorname{Atspat}(\Sigma_{\operatorname{Atspat}(\Sigma)})$
est une bijection.

Pour un atelier quelconque $\mathcal{A}$, de spatios $\Sigma = \Sigma_{\mathcal{A}}$,
le déploiement est un homomorphisme d'ateliers
\[
\mathcal{A} \longrightarrow \operatorname{Atspat}(\Sigma), \qquad X \longmapsto
\operatorname{\text{Dép}}(X),
\]
\emph{fidèle} au sens où $X \mathrel{|\circ|} Y \iff
\operatorname{\text{Dép}}(X) \mathrel{|\circ|} \operatorname{\text{Dép}}(Y)$ ;
et $\mathcal{A}$ est spatial si et seulement s'il est injectif, c'est-à-dire
identifie $\mathcal{A}$ à un sous-atelier de $\operatorname{Atspat}(\Sigma)$ ---
ici, une partie fermée pour $\trianglelefteq$\footnote{Il renvoie à la condition
« Sousat » de sa page 27, notre page 30, et compare à la condition « Atens~1 »
des ateliers ensemblistes, au chapitre VIII.}. Il faut une condition qui tienne
lieu de « tout point est une strate » ; d'où la reformulation sur l'ensemble de
base $\Sigma$ : un spatios, et une partie $\mathcal{A} \subset
\operatorname{Atspat}(\Sigma)$ satisfaisant

\textbf{Atspat 1'.} \emph{$\mathcal{A}$ est fermée pour $\trianglelefteq$ ;}

\textbf{Atspat 2'.} \emph{tout $S \in \Sigma$ est le Sup des $X^{\circ} = \max X$,
$X \in \mathcal{A}$, qui sont $\leq S$.}

\begin{quote}
\textbf{Scholie} (page 59). \emph{La donnée d'un atelier spatial revient à celle
d'un spatios $\Sigma$ et d'une partie $\mathcal{A}$ de
$\operatorname{Atspat}(\Sigma)$ satisfaisant Atspat 1' et Atspat 2' : $\mathcal{A}$
est alors un atelier spatial, $\Sigma \simeq \Sigma_{\mathcal{A}}$, et le
déploiement devient l'inclusion.}
\end{quote}

La vérification (pages 58 et 59) compare $\alpha : X \mapsto \max X$, de
$\mathcal{A}$ dans $\Sigma$, et $\beta : S \mapsto \{X \mid X^{\circ} \leq S\}$,
et conclut : « On gagne ! »\footnote{Elle repose sur ceci : $\alpha$ respecte
exactement la disjonction, et son image est dense pour les Sup (Atspat 2'). C'est
exactement l'énoncé de la proposition 2 des pages 80 à 83, qui caractérise
$\Sigma_{E}$ par ces deux propriétés ; la vérification des pages 58 et 59 en est
un cas particulier écrit avant elle. La page 57 attend que Atspat 1' et 2' soient
« les deux seules (?) » conditions ; c'est ce que la vérification établit pour
l'injectivité et l'identification $\Sigma \simeq \Sigma_{\mathcal{A}}$.}.
L'ensemble de base de la structure, « qui avait été, tour à tour, $\mathfrak{F}$,
$\mathcal{A}$ ou $(\mathcal{A}, \mathfrak{F})$, $\mathcal{L}$, est à présent
$\Sigma$, le “spatios des supports” » --- commode même pour formuler les axiomes
des supports d'ateliers non spatiaux.

\pagerange{60}{74}
\section*{Premier essai sur les correspondances de dispositions (pages 60 à 74)}

Cette section est une première mouture ; la reprise du 10 juillet (pages 75 à
100) la remplace, et c'est elle qu'il faut lire pour les énoncés. On la donne
pour le chemin.

\subsection*{Dispositions et morphismes}

Le « sorite » ouvert au bas de la page 60 pose ceci. Un ensemble $E$ muni d'une relation symétrique $\mathrel{|\circ|}_{E}$ telle que
$x \mathrel{|\circ|} x$ entraîne $x \mathrel{|\circ|} y$ pour tout $y$ est une
\emph{disposition}. Un \emph{morphisme} de $E$ dans $E'$ est une correspondance
$\alpha : E \to \mathfrak{P}(E')$ compatible : si $x \mathrel{|\circ|} y$, tout
élément de $\alpha(x)$ est disjoint de tout élément de $\alpha(y)$. D'où une
catégorie des dispositions\footnote{« (équivalent : un graphe (strict)) » dit la
page 61 ; la correspondance avec les graphes est faite à la page 75.}. On note
$\bar\alpha(A) = \bigcup_{x \in A} \alpha(x)$ et
\[
\alpha_{*} : \Sigma_{E} \to \Sigma_{E'}, \qquad \alpha_{*}(S) =
\operatorname{supp}^{\circ}_{E'}\bigl(\bar\alpha(S)\bigr) .
\]

\subsection*{La formule cherchée, et pourquoi elle manque}

Il voudrait $\alpha_{*}(\operatorname{supp}^{\circ}_{E} A) =
\operatorname{supp}^{\circ}_{E'}(\bar\alpha A)$, et la déduit d'abord d'une
formule « duale » $(**)$, $\alpha_{*}(\operatorname{cosupp}^{\circ} B) =
\operatorname{cosupp}^{\circ}(\bar\alpha B)$, qui revient à dire que
$\alpha_{*}$ commute à $\complement$. Une démonstration directe est essayée et
barrée (« faux ! »). Le contre-exemple est immédiat : pour les dispositions
« triviales », $x \mathrel{|\circ|} y \iff x \neq y$, on a $\Sigma_{E} =
\mathfrak{P}(E)$ ; une application $\alpha : E \to E'$ est compatible si et
seulement si elle est injective, et l'image directe ne commute au complémentaire
que si $\alpha$ est bijective.

Il cherche donc une correspondance en sens inverse $\beta$, propose le choix
« optimal » $\alpha^{*}(x') = \{ x \mid \alpha(x) \subset
\operatorname{supp}^{\circ}\{x'\}\}$, et demande deux conditions --- égalité
$\operatorname{supp}^{\circ}(\bar\alpha(\alpha^{*}(x'))) =
\operatorname{supp}^{\circ}\{x'\}$ et « compatibilité fidèle » $x \mathrel{|\circ|}
y \iff \alpha(x) \mathrel{|\circ|} \alpha(y)$ ---, qui lui paraissent trop
draconiennes. « Donc je n'ai pas l'impression encore d'avoir vraiment bien
compris la situation. »

\subsection*{La transposée}

Il repart (page 67) de l'inclusion qui manque,
$\alpha_{*}(\operatorname{supp}^{\circ} A) \subset
\operatorname{supp}^{\circ}(\bar\alpha A)$, l'autre étant claire. Elle suit de
l'existence d'une correspondance $\beta : E' \to \mathfrak{P}(E)$ telle que
\[
x \mathrel{|\circ|} \beta(x') \iff \alpha(x) \mathrel{|\circ|} x' \qquad (x \in E,
\ x' \in E'),
\]
d'où $A \mathrel{|\circ|} \bar\beta(B') \iff \bar\alpha(A) \mathrel{|\circ|} B'$.
Une telle $\beta$, si elle existe, est déterminée à $\operatorname{supp}^{\circ}$
près : posant $\gamma(x') = \{ x \mid \alpha(x) \mathrel{|\circ|} x'\}$, la
condition dit $\gamma(x') = \operatorname{cosupp}^{\circ}(\beta(x'))$, donc
\emph{$\beta$ existe si et seulement si $\gamma(x') \in \Sigma_{E}$ pour tout
$x'$}, et l'on peut prendre la plus grande, $\beta(x') =
\operatorname{cosupp}^{\circ}(\gamma(x'))$. « Je vois ici déjà poindre une
proposition. »

\begin{quote}
\textbf{Lemme} (page 71). \emph{Une application $\alpha_{*} : \Sigma_{E} \to
\Sigma_{E'}$ qui commute aux Sup quelconques est déterminée par $\alpha =
\alpha_{*} \circ \sigma_{E}$, où $\sigma_{E}(x) = \operatorname{supp}^{\circ}\{x\}$,
et $\alpha_{*}(S) = \operatorname{Sup}_{x \in S} \alpha(x)$.}
\end{quote}

La page énonce une bijection entre les applications $E \to \Sigma_{E'}$ et les
applications de $\Sigma_{E}$ dans $\Sigma_{E'}$ qui commutent aux Sup\footnote{La
bijection est fausse ; seule l'injection vaut. D'abord $\sigma_{E}$ n'est pas
injective en général : si aucun couple n'est disjoint, $\Sigma_{E} = \{\emptyset,
E\}$ et $\sigma_{E}$ est constante, de sorte que seules les $\alpha$ constantes
proviennent d'un $\alpha_{*}$. Ensuite, même quand elle l'est, $S \mapsto
\operatorname{Sup}_{x \in S}\alpha(x)$ ne commute pas aux Sup en général, parce
que le Sup de $S_{i}$ dans $\Sigma_{E}$ est
$\operatorname{supp}^{\circ}(\bigcup S_{i})$, qui peut contenir des éléments
qu'aucun $S_{i}$ ne contient. La marge de la page 71 le pressent (« il n'est pas
clair du tout que \ldots »). L'image de l'injection est décrite par la
proposition 4 de la page 94.}. Le lemme 2 (page 72) donne alors trois conditions
équivalentes pour un tel $\alpha_{*}$ : (i) il existe $\beta_{*}$ commutant aux
Sup avec $\alpha_{*}(S) \mathrel{|\circ|} S' \iff S \mathrel{|\circ|}
\beta_{*}(S')$ ; (ii) il existe $\beta$ comme ci-dessus ; (iii) $\gamma(x') \in
\Sigma_{E}$ pour tout $x'$. $\beta_{*}$ est alors unique.

\textbf{Définition} (page 73). Les $\alpha_{*}$ qui satisfont ces conditions sont
les \emph{correspondances de dispositions}, et $\beta_{*}$ est leur
\emph{transposée} ${}^{t}\alpha_{*}$. La transposition est une bijection
involutive de $\operatorname{Corrdis}(E, E')$ sur $\operatorname{Corrdis}(E',
E)$, et l'on a $\alpha_{*}(\operatorname{supp}^{\circ}_{E} A) =
\operatorname{supp}^{\circ}_{E'}(\bar\alpha A)$. Il ignore si cette dernière
formule, à elle seule, caractérise les correspondances de
dispositions\footnote{La réponse est oui : c'est l'équivalence (i) $\iff$ (iii) de
la proposition 4, page 94.}, et propose d'appeler \emph{application} une
correspondance dont la transposée commute à $\complement$, comme dans le cas
ensembliste où la transposée d'une application est l'image réciproque.

\pagerange{75}{87}
\section*{Reprise du 10 juillet : dispositions et spatios (pages 75 à 87)}

« Je reprends dès le début la formulaire des dispositions. »

\subsection*{1. Dispositions}

Une \emph{disposition} est un ensemble $E$ muni d'une relation
$\mathrel{|\circ|}$ telle que
\begin{enumerate}
\item[Disp 1.] $x \mathrel{|\circ|} y \Rightarrow y \mathrel{|\circ|} x$ ;
\item[Disp 2.] $x \mathrel{|\circ|} x \Rightarrow x \mathrel{|\circ|} y$ pour tout
$y$.
\end{enumerate}
Soit $E_{0}$ l'ensemble des $x$ tels que $x \mathrel{|\circ|} x$ ; la disposition
est \emph{antiréflexive} si $E_{0} = \emptyset$. Les paires non disjointes,
$\{\{x, y\} \mid x \neq y,\ x \text{ et } y \text{ non disjoints}\}$, forment un
graphe simple sur $E$, et l'on récupère la disposition à partir de ce graphe et
de $E_{0}$ : \emph{les dispositions antiréflexives sur $E$ sont exactement les
graphes simples sur $E$}, via le graphe de non-disjonction. La disposition
\emph{triviale}, $x \mathrel{|\circ|} y \iff x \neq y$, correspond au graphe
discret, et $\Sigma_{E} = \mathfrak{P}(E)$.

On note $\Sigma_{E}$ l'ensemble des supports, stable par intersections
quelconques, de plus grand élément $E$ et de plus petit $E_{0}$ ; ses Sup ne sont
pas les réunions en général. $\complement S = \operatorname{cosupp}^{\circ}(S)$
est une involution, et $\sigma_{E}(x) = \operatorname{supp}^{\circ}\{x\}$ donne
une application canonique $\sigma_{E} : E \to \Sigma_{E}$, qui respecte
exactement la disjonction : $x \mathrel{|\circ|} y \iff \sigma(x)
\mathrel{|\circ|} \sigma(y)$. Comme $\operatorname{cosupp}^{\circ}(A) =
\operatorname{cosupp}^{\circ}(A \cup E_{0})$, la restriction $S \mapsto S
\smallsetminus E_{0}$ identifie $\Sigma_{E}$ à $\Sigma_{E \smallsetminus E_{0}}$,
compatiblement à l'ordre et à $\complement$ ; on ne perd rien à supposer la
disposition antiréflexive. Le nom actuel est \emph{espace
d'orthogonalité}\footnote{Un ensemble muni d'une relation symétrique et
irréflexive $\perp$ est un espace d'orthogonalité ; la notion est due à
J.~C.~Dacey (thèse, 1968) et a été développée par D.~J.~Foulis et C.~H.~Randall
dans les années 1970, comme cadre de la logique quantique. Le rapprochement est le
nôtre.}.

\subsection*{2. Spatios}

La définition 1 (page 77) reprend Spat~1 à Spat~3 sous le nom de « spatiale » ;
on en tire les lois de De Morgan, $\complement(\operatorname{Sup} x_{i}) =
\operatorname{Inf}\complement x_{i}$ et inversement, et $\complement 0 = 1$. Un
spatios porte une disposition canonique, $x \mathrel{|\circ|} y \iff x \leq
\complement y$, décroissante en chacun des deux termes. Pour toute disposition
$E$, $(\Sigma_{E}, \subset, \complement)$ est un spatios : Spat~1 et Spat~2 ne
demandent que Disp~1 ; Spat~3 vient de Disp~2, qui donne $A \cap
\operatorname{cosupp}^{\circ}(A) \subset E_{0}$, avec égalité si $A \in
\Sigma_{E}$.

\begin{quote}
\textbf{Proposition 1} (page 79). \emph{Soient $\Sigma$ un spatios et
$\mathrel{|\circ|}$ sa disposition canonique. L'application $\sigma : \Sigma \to
\Sigma_{\Sigma}$, $x \mapsto \operatorname{supp}^{\circ}\{x\} = \{y \mid y \leq
x\}$, est un isomorphisme de spatios, d'inverse $S \mapsto \operatorname{Sup}S$.}

\textbf{Scholie.} \emph{Sur un ensemble donné, les structures de spatios
correspondent bijectivement aux dispositions pour lesquelles l'application
canonique $\Sigma \to \Sigma_{\Sigma}$ est bijective.}
\end{quote}

Les spatios sont donc des dispositions particulières, qu'il propose d'appeler
\emph{canoniques}\footnote{C'est le fait classique que tout treillis
orthocomplémenté complet est le treillis des fermés de son propre espace
d'orthogonalité, $x \perp y \iff x \leq y^{\perp}$. Le rapprochement est le
nôtre.}. Et pour toute disposition $E$, il voudrait voir $\Sigma_{E}$ comme sa
spatios \emph{enveloppe}, caractérisée à isomorphisme unique près :

\begin{quote}
\textbf{Proposition 2} (pages 80 à 83). \emph{Soient $E$ une disposition, $\Sigma$
un spatios et $\alpha : E \to \Sigma$ une application. Pour qu'il existe un
isomorphisme de spatios $\tilde\alpha : \Sigma_{E} \to \Sigma$ tel que
$\tilde\alpha \circ \sigma_{E} = \alpha$, il faut et il suffit que}
\end{quote}
\begin{enumerate}
\item[a)] \emph{$x \mathrel{|\circ|} y \iff \alpha(x) \mathrel{|\circ|} \alpha(y)$
pour tous $x, y \in E$ ;}
\item[b)] \emph{tout $S \in \Sigma$ est le Sup des $\alpha(x)$ qui sont $\leq S$.}
\end{enumerate}
\emph{L'isomorphisme $\tilde\alpha$ est alors unique, donné par $\tilde\alpha(A) =
\operatorname{Sup}_{x \in A}\alpha(x)$, d'inverse $\beta(S) = \{x \mid \alpha(x)
\leq S\}$.}

La condition b) dit que $\alpha(E)$ est \emph{dense pour les Sup} dans $\Sigma$.
La démonstration est celle-ci\footnote{La page écrit $\sigma$ pour la flèche
oblique du diagramme, où l'on attend $\alpha$ ; la suite de la page dit bien
$\alpha$. Elle annonce une formulation équivalente b'), donnée à la page 86.}.
$\tilde\alpha$ est croissante, surjective par b). Pour voir que $\beta(S) \in
\Sigma_{E}$, on montre $\beta(S) = \operatorname{cosupp}^{\circ}(\beta(\complement
S))$ : $x$ est disjoint de tous les $y$ tels que $\alpha(y) \leq \complement S$
si et seulement si $\alpha(x)$ est disjoint de leur Sup, qui est $\complement S$
par b), c'est-à-dire si $\alpha(x) \leq S$. On a $\tilde\alpha\beta = \mathrm{id}$
par b), et $\beta\tilde\alpha = \mathrm{id}$ par le

\begin{quote}
\textbf{Lemme} (page 82). \emph{Dans un spatios, si $A \leq B$, $A' \leq B'$, $B
\mathrel{|\circ|} B'$ et $A' = \complement A$, alors $A = B$ et $A' = B'$}
\end{quote}

--- car $A \leq B \leq \complement B' \leq \complement A' = A$ ---,
appliqué à $A$, $\complement A$ et leurs images $\beta\tilde\alpha(A)$,
$\beta\tilde\alpha(\complement A)$. L'unicité vient de ce que tout $A \in
\Sigma_{E}$ est le Sup des $\sigma_{E}(x)$, $x \in A$. C'est l'analogue exact de
la caractérisation de la complétion de Dedekind--MacNeille d'un ensemble ordonné
par une partie dense pour les Sup et pour les Inf\footnote{H.~M.~MacNeille,
« Partially ordered sets » (1937) ; la caractérisation par la double densité est
de B.~Banaschewski (1956). Ici la densité pour les Inf est fournie par
$\complement$. Le rapprochement est le nôtre.}.

\subsection*{Parties denses}

\begin{quote}
\textbf{Corollaire 1} (pages 83 à 86). \emph{Soient $E$ une disposition et $E'
\subset E$, avec la disposition induite. Conditions équivalentes :}
\end{quote}
\begin{enumerate}
\item[(i)] \emph{$\sigma_{E}(E')$ est dense pour les Sup dans $\Sigma_{E}$ ;}
\item[(ii)] \emph{$S \mapsto S \cap E'$ envoie $\Sigma_{E}$ dans $\Sigma_{E'}$ et
est injective}\footnote{La rédaction de (ii) est inachevée sur la page : le mot
« surjection » est biffé, et la parenthèse ajoute « elle induit donc une iso
$\varphi$ ». Nous donnons la forme que la démonstration utilise, la condition b')
du corollaire 2 appliquée à $E' \to \Sigma_{E}$.} ;
\item[(iii)] \emph{$\psi : \Sigma_{E'} \to \Sigma_{E}$, $S' \mapsto
\operatorname{supp}^{\circ}_{E}(S')$, est surjective ;}
\item[(iv)] \emph{$\sigma(E')$ est dense pour les Sup dans $\sigma(E)$ ;}
\item[(v)] \emph{pour $x, y \in E$, si tous les $x', y' \in E'$ tels que $\sigma(x')
\leq \sigma(x)$ et $\sigma(y') \leq \sigma(y)$ sont disjoints, alors $x
\mathrel{|\circ|} y$}\footnote{La page n'écrit pas où courent $x'$, $y'$ ; si
c'était dans $E$, la condition serait triviale ($x' = x$). L'ordre est l'ordre
canonique de $E$, induit par celui de $\Sigma_{E}$ via $\sigma_{E}$.}.
\end{enumerate}
\emph{Alors $S \mapsto S \cap E'$ et $\psi$ sont des isomorphismes de spatios
inverses l'un de l'autre.}

(i) $\iff$ (iv) par transitivité de la densité --- « dense » au sens des Sup, qui
n'a rien à voir avec la topologie de l'ordre, insiste-t-il. (i) $\Rightarrow$ (v)
parce que dans un spatios, $x = \operatorname{Sup}x_{i}$, $y =
\operatorname{Sup}y_{j}$ et $x_{i} \mathrel{|\circ|} y_{j}$ pour tous $i, j$
entraînent $x \mathrel{|\circ|} y$. (v) $\Rightarrow$ (iv) par le

\begin{quote}
\textbf{Lemme} (page 85). \emph{Soient $\Sigma$ un spatios, $E$ une partie dense
pour les Sup, $S'$ une partie de $\Sigma$ et $\xi \geq \operatorname{Sup}S'$.
Pour que $\xi = \operatorname{Sup}S'$, il faut et il suffit que tout $y \in E$
disjoint de $S'$ soit disjoint de $\xi$.}
\end{quote}

La condition dit que tout $y \in E$ sous $\operatorname{Inf}\complement
\xi'$ est sous $\complement\xi$ ; par densité, $\operatorname{Inf}_{\xi' \in S'}
\complement\xi' \leq \complement\xi$, d'où $\xi \leq \operatorname{Sup}S'$. Enfin
(iii) $\Rightarrow$ (i) est clair, et (ii) $\Rightarrow$ (i) est le

\begin{quote}
\textbf{Corollaire 2} (page 86). \emph{Sous la condition a) de la proposition 2,
la condition b) équivaut à : b') $S \mapsto \beta(S) = \{x \mid \alpha(x) \leq
S\}$ envoie $\Sigma$ dans $\Sigma_{E}$ et est injective.}
\end{quote}

En effet, si b') vaut, $S' = \operatorname{Sup}\alpha(\beta(S)) \leq S$
vérifie $\beta(S') = \beta(S)$, donc $S' = S$.

\begin{quote}
\textbf{Corollaire 5} (page 87). \emph{Soit $E$ une disposition munie d'un ordre
$\leq$ plus fin que l'ordre canonique, et $E' \subset E$ telle que, pour $x, y
\in E$, $E'_{\leq x} \mathrel{|\circ|} E'_{\leq y}$ entraîne $x \mathrel{|\circ|}
y$. Alors $A \mapsto A \cap E'$ et $A' \mapsto \operatorname{supp}^{\circ}_{E}(A')$
sont des isomorphismes inverses entre $\Sigma_{E}$ et $\Sigma_{E'}$}\footnote{La
page dit « une disposition prédonnée $(E, \leq, \mathrel{|\circ|})$ » ; l'hypothèse
que $\leq$ est plus fin que l'ordre canonique (si $x' \leq x$, tout ce qui est
disjoint de $x$ l'est de $x'$) est la nôtre, et c'est ce qui rend l'hypothèse
plus forte que (v), comme dit la page. Les corollaires 3 et 4 n'existent pas :
la numérotation saute de 2 à 5.}.
\end{quote}

C'est ce qu'il fallait à la page 11 : avec pour ordre
$\mathrel{\mathring{\ll}}$, qui est plus fin que l'ordre canonique par At~3,
$\mathcal{A}(L)$ est dense dans $\widehat{\mathcal{A}}(L)$ au sens du corollaire
5 dès que le lemme de la page 11 vaut dans tous les cas, et les supports des deux
ateliers se correspondent.

\pagerange{87}{100}
\section*{Correspondances et transposition (pages 87 à 100)}

La section porte, à la page 87, un titre souligné : « Correspondances entre
dispositions et entre systèmes »\footnote{Précédé d'un numéro lu « 8 bis » avec
doute.}.

\subsection*{L'analogie ensembliste}

Pour deux ensembles, une relation $R \subset E \times E'$ est une application
$\alpha : E \to \mathfrak{P}(E')$, qui se prolonge en $\bar\alpha :
\mathfrak{P}(E) \to \mathfrak{P}(E')$, $A \mapsto \bigcup_{x \in A}\alpha(x)$ ;
$\bar\alpha$ commute aux Sup, et toute application qui commute aux Sup est de
cette forme, puisque $A = \operatorname{Sup}_{x \in A}\{x\}$. La transposée
$\beta$ de $\alpha$ se caractérise, en termes des seules opérations de
$\mathfrak{P}(E)$ et $\mathfrak{P}(E')$, par
\[
\bar\alpha(A) \mathrel{|\circ|} A' \iff A \mathrel{|\circ|} \bar\beta(A'),
\]
puisque les deux membres disent qu'aucun couple $(x, x')$ de $A \times A'$ n'est
dans $R$\footnote{La page écrit dans le membre de droite le signe sans le petit
rond ; c'est la même relation, la disjonction des dispositions triviales.}. C'est
cette caractérisation qui se transporte.

\subsection*{Transposées}

\begin{quote}
\textbf{Proposition 3} (page 89). \emph{Soient $E$, $E'$ deux dispositions,
$\alpha : E \to \mathfrak{P}(E')$ et $\beta : E' \to \mathfrak{P}(E)$. Conditions
équivalentes :}
\end{quote}
\begin{enumerate}
\item[(i)] $x \mathrel{|\circ|} \beta(x') \iff \alpha(x) \mathrel{|\circ|} x'$
\emph{pour $x \in E$, $x' \in E'$ ;}
\item[(ii)] $A \mathrel{|\circ|} \bar\beta(A') \iff \bar\alpha(A)
\mathrel{|\circ|} A'$ \emph{pour $A \subset E$, $A' \subset E'$ ;}
\item[(iii)] \emph{avec $\alpha_{*}(A) = \operatorname{supp}^{\circ}_{E'}(\bar\alpha
A)$ et $\beta_{*}(A') = \operatorname{supp}^{\circ}_{E}(\bar\beta A')$, on a $A
\mathrel{|\circ|} \beta_{*}(A') \iff \alpha_{*}(A) \mathrel{|\circ|} A'$ pour $A
\in \Sigma_{E}$, $A' \in \Sigma_{E'}$.}
\end{enumerate}

On dit alors que $\alpha$ et $\beta$ sont \emph{transposées} l'une de l'autre.
(i) $\iff$ (ii) est immédiat, et (ii) $\iff$ (iii) parce que les deux membres de
(ii) ne changent pas quand on remplace chaque terme par son
$\operatorname{supp}^{\circ}$. Le point technique est le

\begin{quote}
\textbf{Corollaire} (page 91). \emph{Sous ces conditions,
$\operatorname{supp}^{\circ}_{E'}(\bar\alpha(A)) =
\operatorname{supp}^{\circ}_{E'}(\bar\alpha(\operatorname{supp}^{\circ}_{E}A))$
pour toute partie $A$ de $E$, et de même pour $\beta$. En particulier
$\alpha_{*}(\operatorname{supp}^{\circ}_{E}\{x\}) = \alpha'(x)$, où $\alpha'(x) =
\operatorname{supp}^{\circ}_{E'}\alpha(x)$, et $\alpha_{*}(S) =
\operatorname{Sup}_{x \in S}\alpha'(x)$.}
\end{quote}

Par (ii), un $A' \in \Sigma_{E'}$ est disjoint de $\bar\alpha(A)$ si et
seulement si $\beta_{*}(A')$ est disjoint de $A$, donc de
$\operatorname{supp}^{\circ}A$, donc si et seulement si $A'$ est disjoint de
$\bar\alpha(\operatorname{supp}^{\circ}A)$. Les conditions de la proposition ne
dépendent donc que de $\alpha'$ et $\beta'$, et la transposée $\beta_{*}$ est
unique : $\beta_{*}(S')$ est connu dès qu'on connaît les $S$ qui lui sont
disjoints, et ce sont ceux pour lesquels $\alpha_{*}(S) \mathrel{|\circ|} S'$.

\begin{quote}
\textbf{Proposition 4} (page 94). \emph{Soient $\alpha : E \to \mathfrak{P}(E')$,
$\alpha'(x) = \operatorname{supp}^{\circ}_{E'}(\alpha(x))$ et $\alpha_{*}(S) =
\operatorname{supp}^{\circ}_{E'}(\bar\alpha(S)) = \operatorname{Sup}_{x \in
S}\alpha'(x)$. Conditions équivalentes :}
\end{quote}
\begin{enumerate}
\item[(i)] \emph{$\alpha$ admet une transposée ;}
\item[(ii)] \emph{$\alpha_{*}$ commute aux Sup quelconques, et
$\alpha_{*}(\operatorname{supp}^{\circ}_{E}\{x\}) =
\operatorname{supp}^{\circ}_{E'}(\alpha(x))$ pour tout $x$ ;}
\item[(iii)] \emph{$\alpha_{*}(\operatorname{supp}^{\circ}_{E}A) =
\operatorname{supp}^{\circ}_{E'}(\bar\alpha A)$ pour toute partie $A$ de $E$ ;}
\item[(iv)] \emph{pour tout $S' \in \Sigma_{E'}$, $\{x \mid \alpha(x) \subset
S'\} \in \Sigma_{E}$ ;}
\item[(v)] \emph{pour tout $x' \in E'$, $\gamma(x') = \{x \mid \alpha(x)
\mathrel{|\circ|} x'\} \in \Sigma_{E}$.}
\end{enumerate}

La page démontre (i) $\Rightarrow$ (ii) $\Rightarrow$ (iii) $\Rightarrow$ (iv)
$\Rightarrow$ (i). (ii) $\Rightarrow$ (iii) : $\operatorname{supp}^{\circ}A$ est
le Sup des $\operatorname{supp}^{\circ}\{x\}$. (iii) $\Rightarrow$ (iv) : si $S
= \{x \mid \alpha(x) \subset S'\}$, alors $\bar\alpha(S) \subset S'$, donc par
(iii) $\bar\alpha(\operatorname{supp}^{\circ}S) \subset S'$ et
$\operatorname{supp}^{\circ}S \subset S$. (iv) $\Rightarrow$ (i) : on prend $S' =
\operatorname{cosupp}^{\circ}\{x'\}$, qui donne $\gamma(x') \in \Sigma_{E}$, et la
transposée $\beta(x') = \complement\gamma(x')$, comme à la page 69. Quant à (v),
c'est le cas particulier de (iv) qu'on vient d'utiliser, et il l'entraîne,
puisque pour $S' \in \Sigma_{E'}$,
\[
\{x \mid \alpha(x) \subset S'\} = \{x \mid \alpha(x) \mathrel{|\circ|}
\complement S'\} = \bigcap_{x' \in \complement S'} \gamma(x')
\]
est une intersection de supports\footnote{La page établit (v) par un autre
chemin, en passant par une clause b) de (v) qu'elle dit avoir « dû ajouter après
coup » et qui n'est pas dans l'énoncé ; l'argument ci-dessus, le nôtre, montre
qu'elle n'est pas nécessaire.}.

\subsection*{Trois descriptions des correspondances}

\begin{quote}
\textbf{Scholie} (pages 96 à 98). \emph{Soient $\Sigma = \Sigma_{E}$, $\Sigma' =
\Sigma_{E'}$. On a des bijections canoniques}
\[
\operatorname{Corr}(\Sigma, \Sigma') \xrightarrow{\ \varphi\ }
\operatorname{Corr}(E, E') \xrightarrow{\ \psi\ } \operatorname{Corr}^{*}(E, E'),
\]
\emph{où $\operatorname{Corr}(\Sigma, \Sigma')$ est l'ensemble des applications
$\Sigma \to \Sigma'$ qui commutent aux Sup quelconques ;
$\operatorname{Corr}(E, E')$ celui des $\alpha : E \to \Sigma_{E'}$ satisfaisant
la proposition 4 ; $\operatorname{Corr}^{*}(E, E')$ celui des relations $R \subset
E \times E'$ telles que $R(x) \in \Sigma_{E'}$ et ${}^{t}R(x') \in \Sigma_{E}$
pour tous $x$, $x'$. On a $\varphi(\alpha_{*}) = \alpha_{*} \circ \sigma_{E}$ et
$\psi(\alpha) = \{(x, x') \mid \alpha(x) \mathrel{|\circ|} x'\}$.}
\end{quote}

Ainsi $\psi(\alpha)$ n'est pas le graphe de $\alpha$ mais celui de $\complement
\circ \alpha$ --- comme, pour les ensembles, la relation « $x'$ n'est pas dans
$R(x)$ ». La page 98 vérifie que tout $\alpha_{*}$ qui commute aux Sup provient
de $\alpha = \alpha_{*} \circ \sigma_{E}$, qui satisfait la proposition
4\footnote{On le voit aussi directement : une application $f : \Sigma \to
\Sigma'$ entre treillis complets qui commute aux Sup a un adjoint à droite $g$,
$f(S) \leq T \iff S \leq g(T)$, et $f(S) \mathrel{|\circ|} S'$ s'écrit $f(S) \leq
\complement S'$, c'est-à-dire $S \leq g(\complement S')$, soit $S
\mathrel{|\circ|} \complement g(\complement S')$. La transposée est donc
$\complement \circ g \circ \complement$, qui commute aux Sup, et elle existe
toujours. L'argument est le nôtre.}.

\subsection*{Une catégorie autoduale}

Le 11 juillet (?)\footnote{Date marginale lue « 11.7 » avec doute.}, il fait des
dispositions une catégorie : les morphismes $E \to E'$ sont les correspondances
$\alpha : E \to \Sigma_{E'}$ de la proposition 4, la composée de $\alpha$ et de
$\beta : E' \to \Sigma_{E''}$ est $x \mapsto
\operatorname{supp}^{\circ}_{E''}\bigl(\bigcup_{x' \in \alpha(x)}\beta(x')\bigr)$,
et l'identité est $\sigma_{E}$.

\begin{quote}
\emph{La catégorie des dispositions et de leurs correspondances est équivalente à
celle des spatios et des applications qui commutent aux Sup quelconques. Cette
catégorie est équivalente à son opposée, par une autodualité qui est l'identité
sur les objets et la transposition sur les morphismes.}
\end{quote}

« Analogie avec les espaces de Hilbert, plus généralement espaces vectoriels munis
d'une autodualité, et les opérateurs adjoints. » En termes actuels : la catégorie
des treillis orthocomplémentés complets et des applications qui préservent les
Sup est une \emph{catégorie à poignard} (dagger category), le poignard étant la
transposition\footnote{B.~Jacobs, « Orthomodular lattices, Foulis semigroups and
dagger kernel categories » (2010), étudie précisément ces catégories, dont
l'exemple de référence est celui des espaces de Hilbert et des adjoints. Le
rapprochement est le nôtre ; il n'y a aucune raison de penser que ces feuillets
aient été connus de ces auteurs.}.

\subsection*{Questions ouvertes}

Sous ce titre (page 99), qu'il n'a pas « tirées au clair » : pour des ensembles,
les correspondances forment elles-mêmes l'ensemble $\mathfrak{P}(E \times E')$,
système associé à la disposition triviale $E \times E'$. Il serait tentant de voir
les correspondances entre deux dispositions comme les éléments du spatios de la
\emph{disposition produit}
\[
(x, x') \mathrel{|\circ|} (y, y') \iff x \mathrel{|\circ|} y \ \text{ou}\ x'
\mathrel{|\circ|} y' .
\]
Les éléments\footnote{La page écrit « $x \mathrel{|\circ|} x'$ ou $y
\mathrel{|\circ|} y'$ », qui compare des éléments d'ensembles différents ; la
lecture ci-dessus est celle que le transcripteur indique et la seule qui ait un
sens. Pour les dispositions triviales elle redonne bien « $(x,x') \neq (y,y')$ ».} de $\Sigma_{E \times E'}$ sont bien dans
$\operatorname{Corr}^{*}(E, E')$, mais la réciproque ne lui semble vraie que si
l'une des dispositions est triviale ; et la question demeure de savoir si les
correspondances forment un spatios. C'est, sous une autre forme, la question du
produit tensoriel que la logique quantique a rencontrée dans les mêmes
années\footnote{D.~J.~Foulis et C.~H.~Randall, « Empirical logic and tensor
products » (1981) : le produit tensoriel de deux treillis orthomodulaires
n'existe pas en général. Le rapprochement est le nôtre.}.

\pagerange{101}{107}
\section*{L'analogie ensembliste et ses questions (pages 101 à 107)}

Sur $\operatorname{Corr}(\Sigma, \Sigma')$, l'ordre $f \leq g$ si $f(x) \leq g(x)$
pour tout $x$ a des Sup quelconques, calculés point par point. Mais la structure
$\complement$, ou de disposition, n'est pas claire. Pour $E$ trivial,
$\operatorname{Corr}(\Sigma_{E}, \Sigma') \simeq \Sigma'^{E}$, avec $\complement$ et
$\mathrel{|\circ|}$ terme à terme --- mais seulement sur les points : sur les
parties, $(\complement f)(A) \neq \complement(f(A))$ en général. La question lui
semble liée à l'existence d'éléments minimaux de $\Sigma \smallsetminus \{0\}$
(les atomes), « ou d'autres objets qui en tiendraient lieu, jouant le rôle des
“points” ».

\subsection*{Comorphismes et morphismes}

Pour des ensembles, une correspondance de graphe $\Gamma$, $E \xleftarrow{f}
\Gamma \xrightarrow{f'} E'$, se décompose en $\alpha_{*}(A) = f'(f^{-1}(A))$ : une
image réciproque (« comorphisme ») suivie d'une image directe (« morphisme »). Un
comorphisme se caractérise par la commutation à $\complement$, un morphisme par
celle de sa transposée. Y a-t-il une décomposition semblable pour les
dispositions ? Ce n'est pas, dit-il, de conséquence directe pour la
formalisation ensembliste des ateliers ; les questions suivantes le sont.

\begin{quote}
\textbf{A)} (page 104). \emph{Pour $\alpha_{*} : \mathfrak{P}(E) \to
\mathfrak{P}(E')$ qui commute aux Sup, sont équivalents : $\alpha_{*}$ respecte
la disjonction ; $\alpha_{*}$ commute aux Inf non vides ; $\alpha_{*}$ commute à
$S \wedge \complement T$ (ou seulement pour $T \leq S$) ; $f'$ est injective.}

\textbf{B)} \emph{Sont équivalents : $\alpha_{*}(1) = 1$ ; $\beta_{*}(S) = 0
\Rightarrow S = 0$ ; $f'$ est surjective.}
\end{quote}

D'où trois corollaires. \emph{Corollaire 1} : $\alpha_{*}$ respecte la
disjonction et $\alpha_{*}(1) = 1$ si et seulement s'il commute aux Inf
quelconques, si et seulement s'il commute à $\complement$, si et seulement si
$\alpha$ est un comorphisme $A \mapsto g^{-1}(A)$, $g : E' \to E$\footnote{La page
écrit « commute à $\complement$, i.e.\ $f : \Gamma \xrightarrow{\sim} E$ » et $g =
f' f^{-1}$. Avec ses conventions ($f$ vers $E$, $f'$ vers $E'$), c'est $f'$ qui
doit être bijective, par A) et B) ; et $g = f \circ f'^{-1}$, comme le
transcripteur le relève.}. \emph{Corollaire 2} : $\alpha_{*}$ respecte la
disjonction « fidèlement », $S \mathrel{|\circ|} T \iff \alpha_{*}(S)
\mathrel{|\circ|} \alpha_{*}(T)$, si et seulement s'il la respecte et ne tue
aucune partie non vide, si et seulement si $\alpha$ est un comorphisme le long
d'une surjection suivi d'un morphisme injectif. \emph{Corollaire 3} : commuter à
$\complement$ sans tuer de partie non vide, ou respecter fidèlement la
disjonction avec $\alpha_{*}(1) = 1$, c'est être un comorphisme le long d'une
surjection $g : E' \to E$.

\subsection*{Le cas général}

Ces équivalences gardent-elles un sens pour des dispositions quelconques ? Sans
décomposition canonique $\alpha_{*} = f'_{*}f^{*}$, l'interprétation directe
disparaît ; il reste des équivalents algébriques, dont les preuves dans A)
passent par des hypothèses « draconiennes », booléennes. Mais ce sont peut-être
de faux problèmes : pour le travail géométrique, on n'aura affaire qu'aux
supports constructibles, et aux Sup et Inf de familles « commensurables ». Dans
l'intuition géométrique, respecter la disjonction, c'est être un comorphisme
suivi d'une injection ; avec $\alpha_{*}(1) = 1$, on s'attend à la commutation à
$\complement$.

Deux remarques valent dans tout spatios : la commutation à $\complement$ entraîne
$\alpha_{*}(1) = 1$ et le respect de la disjonction ; et l'équivalence des deux
premières conditions de B) « se vérifie trivialement »\footnote{En effet
$\beta_{*}(S) = 0$ équivaut à $1 \mathrel{|\circ|} \beta_{*}(S)$, donc à
$\alpha_{*}(1) \mathrel{|\circ|} S$, c'est-à-dire $S \leq
\complement\alpha_{*}(1)$.}. Il pose enfin la question : si $\alpha_{*}$ et sa
transposée commutent toutes deux à $\complement$, sont-elles des isomorphismes de
spatios inverses l'un de l'autre ? La réponse est oui\footnote{La question est lue
en partie avec doute. Si la transposée $\beta_{*}$ commute à $\complement$, elle
est égale à l'adjoint à droite $g$ de $\alpha_{*}$ (puisqu'elle vaut $\complement
g \complement$), donc $\alpha_{*}\beta_{*} \leq \mathrm{id} \leq
\beta_{*}\alpha_{*}$. Et $\alpha_{*}\beta_{*}(\complement S) =
\complement\alpha_{*}\beta_{*}(S) \leq \complement S$ donne $\alpha_{*}\beta_{*}(S)
\geq S$, d'où $\alpha_{*}\beta_{*} = \mathrm{id}$ ; de même $\beta_{*}\alpha_{*} =
\mathrm{id}$. L'argument est le nôtre.}.

\pagerange{108}{121}
\section*{Orthomodularité, familles admissibles, axiomes des supports (pages 108 à 121)}

\subsection*{Deux éléments comparables}

\begin{quote}
\textbf{Lemme 1} (page 108). \emph{Soient $b \leq c$ dans un spatios. Il existe
$b^{*}$ disjoint de $b$ tel que $b \vee b^{*} = c$ si et seulement si $b \vee
(\complement b \wedge c) = c$ ; et $b^{*}_{0} = \complement b \wedge c$ est le plus
grand de ces $b^{*}$.}

\textbf{Lemme 2.} \emph{Pour $b \leq c$, sont équivalents : (i) $b$ et $c$
engendrent une sous-algèbre de Boole ; (ii) ils engendrent une algèbre de Boole
finie, de cardinal $2^{n}$, $n \in \{1, 2, 3\}$ ; (iii) $(b, c)$ et $(\complement
c, \complement b)$ satisfont la condition du lemme 1 ; (iv) $\alpha = b$, $\beta =
c \wedge \complement b$ et $\gamma = \complement c$ sont deux à deux disjoints et
chacun est le complément de la réunion des deux autres.}

\textbf{Corollaire.} \emph{Sont équivalents, pour un spatios $\Sigma$ : (i) tout
couple $b \leq c$ satisfait le lemme 1 ; (ii) tout couple $b \leq c$ satisfait le
lemme 2 ; (iii) deux éléments disjoints engendrent toujours une sous-algèbre de
Boole ; (iiii) de même pour toute famille finie d'éléments deux à deux
disjoints.}
\end{quote}

On a $n = 1$ si $\{b, c\} \subset \{0, 1\}$, $n = 3$ si $0 < b < c < 1$, $n = 2$
sinon ; dans ce dernier cas, $\alpha$, $\beta$, $\gamma$ forment une base de
l'algèbre engendrée (corollaire 2, page 110)\footnote{La dernière ligne de (iv),
page 108, est écrite de façon fautive (un $\wedge$ manque, et les lettres ne
s'accordent pas avec les relations qu'elle reformule) ; nous donnons les relations
$\complement(\alpha \vee \beta) = \gamma$, etc., qu'elle annonce.}. La condition
(i) du corollaire est la \emph{loi orthomodulaire} : $b \leq c \Rightarrow c = b
\vee (b^{\perp} \wedge c)$. Les spatios qui la satisfont sont les \emph{treillis
orthomodulaires} complets, et l'équivalence avec (iii) et (iiii) est la
caractérisation classique par les éléments qui « commutent »\footnote{La loi
orthomodulaire est due à K.~Husimi (1937) ; elle caractérise, parmi les treillis
orthocomplémentés, ceux où deux éléments orthogonaux engendrent une sous-algèbre
de Boole (voir G.~Kalmbach, \emph{Orthomodular Lattices}, 1983). Le treillis des
sous-espaces fermés d'un espace de Hilbert est orthomodulaire, et non distributif
dès la dimension 2. Le nom est le nôtre.}.

\textbf{Question} (page 110). Ces conditions sont-elles toujours vérifiées ?
« J'en doute fort ! » Le sont-elles si $\Sigma$ est fini ? « J'en doute
également. » Si elles le sont, $\Sigma$ est-il une algèbre de Boole ? Les deux
doutes sont fondés, et la dernière réponse est non\footnote{L'hexagone $O_{6} =
\{0, a, b, \complement a, \complement b, 1\}$, avec $a < \complement b$ (donc $b
< \complement a$) et aucune autre relation, est un spatios fini où $a \leq
\complement b$ et $a \vee (\complement a \wedge \complement b) = a \neq
\complement b$ : il n'est pas orthomodulaire. Le treillis $MO_{2} = \{0, a,
\complement a, b, \complement b, 1\}$, où les quatre éléments du milieu sont deux
à deux incomparables, est orthomodulaire sans être une algèbre de Boole. Par la
proposition 1 de la page 79, l'un et l'autre sont le spatios de leur propre
disposition. Les deux exemples sont classiques ; la réponse est la nôtre.}. Il
suggère de chercher un exemple parmi les spatios construits sur un ensemble
ordonné $L$ avec la disposition « être comparables », et la page 112 en dessine
un : un treillis en forme de prisme à section hexagonale, avec des points $i$,
$i'$, $j$, $k$, trois éléments $a'$, $b$, $c$ et leurs réunions $A = b \vee c$,
$B = c \vee a'$, $C = a' \vee b$ ; il vérifie les compléments, trouve « OK » pour
deux d'entre eux et « pas OK ! » pour le troisième, et en reste là\footnote{La
marge de la page 110 porte « contrexemple \ldots\ $\Sigma = 10$ », d'une lecture
incertaine. Le dessin de la page 112 n'est pas assez lisible pour qu'on en
reconstitue la construction ; nous n'en disons que ce que la transcription en
décrit.}.

\subsection*{Familles admissibles}

« Après cette longue digression, il faudrait revenir à l'axiome des supports dans
un atelier. »

\begin{quote}
\textbf{Lemme 1} (page 114). \emph{Soit $(e_{i})_{i \in I}$ une famille
d'éléments non nuls deux à deux disjoints d'un spatios, de Sup égal à $1$. Sont
équivalents : (i) les $e_{i}$ sont contenus dans une partie de $\Sigma$ stable
par Sup, Inf et $\complement$ qui est une algèbre de Boole ; (ii) l'application
$\mathfrak{P}(I) \to \Sigma$, $J \mapsto e_{J} = \operatorname{Sup}_{i \in
J}e_{i}$, qui commute aux Sup, commute à $\complement$. Son image est alors la plus
petite partie satisfaisant (i), et elle est injective.}

\textbf{Corollaire 1} (page 115). \emph{Si $e = \operatorname{Sup}e_{i}$ n'est pas
$1$, on pose $e_{0} = \complement e$ ; la famille $(e_{i}) \cup \{e_{0}\}$
satisfait le lemme si et seulement si $e_{I'} = e \wedge \complement e_{I''}$
pour toute partition $I = I' \sqcup I''$.}
\end{quote}

\textbf{Définition.} Une famille d'éléments non nuls deux à deux disjoints qui
satisfait ces conditions est \emph{admissible} ; elle engendre une algèbre de
Boole $\Sigma_{(e_{i})}$ isomorphe à $\mathfrak{P}(I \sqcup \{0\})$, ou à
$\mathfrak{P}(I)$ si $e_{0} = 0$\footnote{La page écrit seulement
$\mathfrak{P}(I \sqcup \{e_{0}\})$ ; quand $e = 1$, l'élément $e_{0}$ est nul et
ne compte pas. (i) $\Rightarrow$ (ii) utilise la distributivité infinie des
algèbres de Boole complètes ; l'injectivité, notée en marge, vient de ce que $J =
\{i \mid e_{i} \leq e_{J}\}$ sous (ii).}. Il dira plus loin plutôt
« commensurable ».

\textbf{Axiome Atmagsupp 1} (page 116). \emph{Pour toute figure $F$ du magasin,
la famille $(X^{\circ})_{X \in F}$ des supports de ses strates est admissible dans
$\Sigma_{\mathcal{A}}$}\footnote{La page écrit $F \in \mathcal{F}_{c}$, sans doute
le magasin des figures compactes ; le nom de l'axiome est lu « Atmagsupp » avec
doute.}.

L'axiome n'est pas redondant : pour l'atelier
$\operatorname{\text{Figélspat}}_{\mathrm{fin}}(\Sigma)$ des figures spatiales
élémentaires finies d'un spatios, avec le magasin de ses figures spatiales finies,
il équivaut à : si $e \mathrel{|\circ|} e'$, alors $e' = \complement e \wedge (e
\vee e')$ --- c'est-à-dire aux conditions du lemme 2 de la page 108. Autrement
dit, \emph{Atmagsupp 1 vaut pour l'atelier universel d'un spatios si et seulement
si ce spatios est orthomodulaire}\footnote{Avec $e = \complement c$, $e' = b$ pour
$b \leq c$, la condition s'écrit $b = c \wedge (\complement c \vee b)$, qui est la
loi orthomodulaire pour $\complement c \leq \complement b$ ; la réciproque se
démontre dans tout treillis orthomodulaire. La formulation en ces termes est la
nôtre.} ; sur l'hexagone il est faux.

\subsection*{Deux figures, puis un drapeau}

Pour deux parties $A$, $B$ formées d'éléments non nuls deux à deux disjoints ---
des « figures spatiales discrètes » ---, on écrit $B \ll A$ si tout élément de $B$
est sous un élément de $A$. Pour $a \in A$, soient $B_{a} = \{b \in B \mid b \leq
a\}$, $a' = \operatorname{Sup}B_{a}$ et $a'' = a \wedge \complement a'$, ce qui
reste de $a$ hors de $B$.

\begin{quote}
\textbf{Lemme 2} (page 118). \emph{$A \cup B$ est contenu dans une sous-algèbre de
Boole stable par Sup, Inf et $\complement$ si et seulement si a) $a = a' \vee
a''$ pour tout $a \in A$, et b) $\widehat{B} = B \cup \{a'' \mid a'' \neq 0\}$ est
admissible. L'algèbre de Boole engendrée par $\widehat{B}$, isomorphe à
$\mathfrak{P}(\widehat{B})$, est alors la plus petite possible.}
\end{quote}

On dit alors que $\{A, B\}$ est admissible, et l'on pose l'\textbf{axiome
Atmagsupp 2} : \emph{pour $F \ll G$ dans le magasin, avec $B = \{X^{\circ} \mid X
\in F\}$ et $A = \{X^{\circ} \mid X \in G\}$, la paire $\{A, B\}$ est
admissible}\footnote{La page indexe $B$ par $G$ et $A$ par $F$ ; avec $F \ll G$ et
la convention $B \ll A$ de la page 117, c'est l'inverse.}. Mais ce n'est pas
assez : il faut traiter un drapeau quelconque $F_{0} \ll F_{1} \ll \cdots \ll
F_{n}$, qui donne un drapeau $A_{0} \ll \cdots \ll A_{n}$ de figures spatiales
discrètes de $\Sigma_{\mathcal{A}}$ (en oubliant l'ordre de chaque figure).

\textbf{Définition} (pages 119 et 120). Une partie de $\Sigma$ est
\emph{commensurable} si elle est contenue dans un sous-spatios --- stable par Sup,
Inf et $\complement$ --- qui est une algèbre de Boole ; c'est plus fort que la
commensurabilité deux à deux. Un drapeau est commensurable si la réunion de ses
termes l'est\footnote{La page écrit $\bigcup_{1 \leq i \leq n}A_{i}$ ; le lemme 3
qui suit pose $A''_{0} = A_{0}$, et c'est $0 \leq i \leq n$ qu'il faut.}.

\textbf{Lemme 3} (pages 120 et 121). Soient $A''_{i}$ l'ensemble des restes non
nuls $a''_{i}$, $a_{i} \in A_{i}$, définis comme au lemme 2 pour $A_{i-1} \ll
A_{i}$, $A''_{0} = A_{0}$, et $A = \bigcup_{i} A''_{i}$. Les $A''_{i}$ sont
formés d'éléments non nuls deux à deux disjoints, et sont disjoints entre eux ;
pour que le drapeau soit commensurable, il faut et il suffit que $A$ le
soit\footnote{L'énoncé s'interrompt au bas de la page 120 et reprend au haut de
la page 121, en partie illisible ; la conclusion est celle que la page 121 lit
(« il faut et il suffit que $A$ le soit »).}.

\textbf{Axiome Atmagsupp 3} (page 121), qui englobe les deux précédents :
\emph{pour tout drapeau $F_{0} \ll \cdots \ll F_{n}$ de figures, le drapeau
correspondant $A_{0} \ll \cdots \ll A_{n}$ de
$\operatorname{Figspatdisc}(\Sigma_{\mathcal{A}})$ est commensurable.}

Il envisage alors un atelier $\operatorname{\text{Figélspatcomm}}(\Sigma)$ des
figures spatiales élémentaires commensurables, avec une relation
$\mathrel{\mathring{\ll}}_{\mathrm{comm}}$ plus forte. Mais il n'est pas clair
qu'elle soit transitive : si $e = f \vee f'$ et $f = g \vee g'$, avec
$\{g, g'\}$ et $\{g \vee g', f'\}$ commensurables, il ne s'ensuit pas que $\{g,
g', f'\}$ le soit. « Donc ça ne semble pas marcher \ldots\ Ça marche bien sûr
quand $\Sigma$ lui-même est une algèbre de Boole. »\footnote{Ça marche déjà quand
$\Sigma$ est orthomodulaire : $g$, $g'$, $f'$ sont deux à deux disjoints, donc
commutent deux à deux, et une famille finie d'éléments qui commutent deux à deux
engendre une sous-algèbre de Boole (théorème de Foulis et Holland, 1962--1963 ;
voir Kalmbach, \emph{op.\ cit.}). La difficulté qu'il rencontre est celle des
spatios non orthomodulaires. La relation entre $f$ et $f'$ est lue avec doute.}

\pagerange{121}{127}
\section*{Critique philosophique et programme (pages 121 à 127)}

Il faudrait voir si l'existence d'une « fonction support fidèle » au sens du
chapitre VIII entraîne toujours Atmagsupp 3, et regarder le cas des ateliers
« magmatiques » fidèles\footnote{Les renvois sont au chapitre VIII, pages 62 et
78 de sa main ; le mot lu « magmas » ou « magmatiques », ici et aux pages 124 à
129, n'est défini sur aucune page de ce dossier.}. « Mais avant tout, il me
faudrait aller de l'avant, pour me rendre compte exactement comment vont
intervenir les axiomes des supports, et s'ils s'avèrent aussi cruciaux que je le
pressens. »

Suit une section intitulée « Critique philosophique de cette
approche »\footnote{Les pages 123 à 126 sont écrites vite et la transcription y
est très lacunaire ; on n'en rend que ce qui se lit.}. Pour définir les
opérations sur les supports, on est forcé de « couper » des strates qui ne sont
ni compatibles, ni même commensurables. Or, en fin de compte, il ne s'est
intéressé aux supports et aux opérations sur eux que pour les supports
constructibles, et pour des familles commensurables de tels supports. « Il y a là
une contradiction. » Il lui paraît pourtant indispensable d'avoir une description
« intrinsèque » des supports subordonnés à une figure, indépendante de celle-ci,
et de leurs relations avec ceux d'une subdivision, qui ne soit pas affectée par
le fait que deux subdivisions ne soient pas commensurables. Il esquisse deux
voies : un ensemble $\Sigma_{\mathrm{ens}}$ des supports constructibles, avec des
opérations partout définies et d'autres partiellement définies ; ou les supports
constructibles fermés obtenus en « inversant » les flèches $\ll$ de
$(\mathcal{F}, \ll)$, les subdivisions --- une catégorie de fractions dont il
n'est pas sûr qu'elle soit un ensemble ordonné, ni qu'elle se plonge dans
$\Sigma$.

Il se rassure : quand $\Sigma$ est booléen --- ateliers pseudo-booléens, peut-être
les ateliers fidèles, en tout cas les ateliers magmatiques ---, il n'y a « vraiment
rien à lui reprocher », et cela inclut déjà tous les cas qui importent. Ce sont
les « ordres », les pseudo-simpliciaux topologiques, « associés à l'intuition du
“lieu” remplaçant le “point” », qui causent tous les ennuis. Mais même dans le cas
booléen, il serait intéressant de décrire $\Sigma_{\mathrm{ens}}$ de façon
combinatoire, « dans un esprit plus proche d'une “géométrie des formes” --- dans
laquelle la spatialité a de forts relents de topologie générale ! »

Programme : « boucler » l'aspect provisoirement ensembliste --- « il serait
temps ! » --- en traitant a) les subdivisions, « enfin ! », et les morphismes ; b)
les ateliers modérés ; c) les axiomes de saturation.

\pagerange{129}{139}
\section*{IX bis : morphismes d'ateliers, excision, subdivisions (pages 129 à 139)}

Une nouvelle pagination commence, sous l'en-tête « cf IX bis » et un renvoi au
chapitre VIII.

\subsection*{Trois situations}

Il veut une notion de morphisme qui couvre trois cas.

\emph{1° Un sous-atelier} $\mathcal{A}' \subset \mathcal{A}$, au sens de la
condition Sous-at de la page 30 ; muni des relations induites c'est un atelier, et
$\ll_{\mathcal{A}'}$ est induite par $\ll_{\mathcal{A}}$\footnote{Il renvoie à
« X p.~27 », d'une lecture douteuse : la page 27 de sa main est notre page 30.}.
Exemples : l'ombre d'une préfigure (page 30) ; une figure elle-même, qui est un
sous-atelier, même « magmatique », car $X \mathrel{\mathring{\ll}} Y'
\trianglelefteq Y$ avec $X, Y \in F$ force $X = Y'$ ; l'« atelier résiduel
strict » des strates dont l'adhérence est disjointe d'une partie $S$, fermé pour
$\ll$. Et l'\emph{atelier résiduel au sens large} d'une préfigure $\mathcal{F}$,
\[
\mathcal{A}' = \{ X \in \mathcal{A} \mid X \mathrel{|\circ|} \mathcal{F} \text{ et
} X \text{ compatible avec chaque strate de } \mathcal{F} \},
\]
est un sous-atelier\footnote{La page note la compatibilité par un signe que la
transcription rend par $\between$, de forme douteuse ; nous y lisons la relation
$\asymp$ de la page 6, puisque la même page définit une préfigure par « $X, Y \in
\mathcal{F} \Rightarrow X \between Y$ ». La page appelle
$\operatorname{Supp}^{\circ}(S)$ l'ensemble $\{X \mid X \mathrel{|\circ|} S\}$,
notre $\operatorname{cosupp}^{\circ}(S)$, et renvoie pour le contre-exemple à sa
page 28, nos pages 31 et 32.}. En effet, si $X \mathrel{\mathring{\ll}} Y'
\trianglelefteq Y$ avec $X, Y \in \mathcal{A}'$, $Y'$ est compatible avec chaque
$Z \in \mathcal{F}$ parce que $Y$ l'est ; pour deux strates compatibles, être
disjointes revient à être distinctes ; et si l'on avait $Y' = Z$, on aurait $X
\mathrel{\mathring{\ll}} Z$ et $X \mathrel{|\circ|} Z$, donc $X \mathrel{|\circ|}
X$. C'est la version juste de l'affirmation de la page 32 : la compatibilité
exclut le triangle qui la mettait en défaut. Il faudrait voir si $\Sigma_{\mathcal{A}'}$
se plonge dans $\Sigma_{\mathcal{A}}$.

\emph{2° Une application d'ensembles} $f : E' \to E$ (ou d'espaces modérés), et
l'image réciproque des figures. Si $f$ n'est pas surjective, elle envoie une
strate de $\mathcal{A}$ dans $\mathcal{A}' \cup \{\emptyset\}$. Elle respecte
$\trianglelefteq$, $\mathrel{\mathring{\ll}}$, $\mathrel{|\circ|}$, n'est pas
injective, et sur les supports commute à tout : Sup, Inf, $\complement$.

\emph{3° Une injection} $g : E \to E'$ et l'image directe, injective ; elle
respecte les trois relations, et sur les supports commute aux Sup, aux Inf non
vides et à $S \wedge \complement T$, mais pas à $\complement$ si $g$ n'est pas
bijective.

Le cas commun aux deux derniers est une correspondance $E \xleftarrow{f} \Gamma
\xrightarrow{f'} E'$ avec $f'$ injective --- exactement la condition A) de la
page 104. Il ne semble pas y avoir d'application directe des figures pour une $g$
non injective.

\subsection*{Homomorphismes et excision}

\textbf{Définition} (provisoire, page 133). Un \emph{homomorphisme} d'ateliers est
une application $f : \mathcal{A} \to \widehat{\mathcal{A}}' = \mathcal{A}' \sqcup
\{\emptyset\}$ qui respecte $\trianglelefteq$, $\mathrel{\mathring{\ll}}$ et
$\mathrel{|\circ|}$, avec les conventions : $\emptyset \trianglelefteq X'$ et $X'
\trianglelefteq \emptyset$ toujours (ce n'est plus un ordre) ; $\emptyset
\mathrel{\mathring{\ll}} X'$ toujours et $X' \mathrel{\mathring{\ll}} \emptyset$
seulement pour $X' = \emptyset$ (un ordre de plus petit élément $\emptyset$) ;
$X' \mathrel{|\circ|} \emptyset$ toujours\footnote{La convention pour
$\trianglelefteq$ est lue ainsi d'après la remarque que ce n'est pas une relation
d'ordre ; la page écrit « $X' \trianglelefteq \emptyset$, ($\emptyset
\trianglelefteq X'$ toujours vrai) », d'une rédaction biffée en partie. Les
conventions sont inspirées de l'exemple 2°.}. La composition se lit sur les
applications $\mathcal{A} \sqcup \{\emptyset\} \to \mathcal{A}' \sqcup
\{\emptyset\}$ qui fixent $\emptyset$.

Le \emph{noyau} $\operatorname{Ker}f = \{X \mid f(X) = \emptyset\}$ est fermé vers
le bas pour $\mathrel{\mathring{\ll}}$. On s'attend à ce que ce soit un support,
moyennant des conditions supplémentaires (c'est le cas dans l'exemple 2°).
Inversement, une partie $N$ fermée pour $\mathrel{\mathring{\ll}}$ est-elle le
noyau d'un homomorphisme \emph{universel} $p : \mathcal{A} \to \mathcal{A}/N$, par
lequel tout homomorphisme de noyau contenant $N$ se factorise de façon unique ---
une opération d'\emph{excision} ? L'exemple type est $\mathcal{A} =
\operatorname{Figelann}(E)$, $N$ les strates de support contenu dans $E
\smallsetminus E'$, et $\mathcal{A}/N = \operatorname{Figelann}(E')$\footnote{Le
nom $\operatorname{Figelann}$ est lu tel qu'il est écrit, sans certitude ; la fin
de la page 135 suggère une description comme quotient de $\mathcal{A} \smallsetminus
N$, d'une lecture très incertaine.}. La construction « s'annonce délicate » ; il
faudrait regarder d'abord les ateliers ensemblistes --- pas évident, même pour $N
= \operatorname{Omb}^{\circ}F$ --- puis le cas d'un ensemble ordonné, où
$\mathrel{\mathring{\ll}}$ est l'égalité et où l'on doit trouver simplement $A
\smallsetminus N$. « Là je n'ai pas le [temps ?] pour, et laisse tomber pour le
moment. »

\subsection*{Subdivisions}

Soient $G \ll F$ deux figures, et $\varphi : \widetilde{G} \to \widetilde{F}$
l'application croissante telle que $Y \mathrel{\mathring{\ll}} \varphi(Y)$. Pour
que $G$ soit une \emph{subdivision} de $F$, il exige
\begin{enumerate}
\item[(i)] pour tout $X \in \widetilde{F}$, $\varphi^{-1}(X)
\mathrel{\mathring{\ll}} X$ et $\operatorname{supp}^{\circ}\varphi^{-1}(X) =
X^{\circ}$ : les morceaux de $G$ au-dessus de $X$ remplissent $X$ ;
\end{enumerate}
et note que, moyennant l'axiome des supports, cela équivaut à
\begin{enumerate}
\item[(ii)] $\operatorname{supp}G = \operatorname{supp}F$ ;
\item[(iii)] il n'existe pas de strate $Z \ll F$ disjointe de $G$ ;
\item[(iv)] si l'atelier est divisible : il n'existe pas de point $x \ll F$
disjoint de $G$ ;
\item[(v)] si l'atelier est divisible : $G$ est maximal parmi les raffinements de
$F$.
\end{enumerate}
Les implications (i) $\Rightarrow$ (ii) $\Rightarrow$ (iii) $\Rightarrow$ (iv)
sont triviales, et (iv) $\Rightarrow$ (i) par la divisibilité ; reste (iii)
$\Rightarrow$ (i) sans divisibilité, « qu'en Atmagsupp ? »\footnote{La chaîne
d'implications a plusieurs numéros récrits. Le sens de « maximal parmi les
raffinements de $F$, admis par $\mathcal{C}$ » dans (v) n'est pas précisé sur la
page.}. Une condition supplémentaire est essayée puis barrée : que l'ordre de
$\widetilde{F}$ soit le quotient de celui de $\widetilde{G}$ --- que $F$ se
déduise de $G$ par simple « regroupement de strates ». La rature s'arrête au
milieu de la page 139, la dernière du dossier.

\end{document}
